\chapter{外场散射与有限脉冲}\label{chap:external-pinem}

一个电子穿过受激光照射的纳米结构时，近场可以使它获得或失去能量。若光主要含有角频率 $\omega$，出射能谱会在入射中心能量 $E_0$ 两侧出现间距约为 $\hbar\omega$ 的峰。这类利用光学近场调制并读取电子的实验称为\emph{光子诱导近场电子显微}（PINEM）。用整数 $\ell$ 给这些能量边带编号，$\ell>0$ 表示电子净获得能量，$\ell<0$ 表示净损失能量。\chapref{chap:electron-finite-time}已经建立有限时间内的电子--环境跃迁理论，现在用它来求各条边带的振幅和相位。

先考虑光场可由给定经典电磁场描述的情形。此时需要求电子在外场中的波函数，近场的空间分布和光学相位都进入散射振幅。随后再保留光场的产生、湮灭算符，考察光子数涨落以及电子与光的纠缠；这部分通常称为\emph{量子光子诱导近场电子显微}（QPINEM）。两种描述的联系由相干光的经典极限给出，而电子反冲和有限脉冲的影响需要在各自模型中检查。\cite{Feist2015Nature,Wang2020Cavity,Dahan2021Science,Ruimy2025Perspective}

\section{外场散射与相位约化}

给定入射波包和外场四势 $A^\mu(x)$，外场散射要求求出电子离开相互作用区时的振幅。我们从 Dirac 方程的传播子开始。当入射动量集中在中心值附近、外场也不明显改变电子前进方向时，可以把波函数分成快速振荡的载波与缓慢变化的包络。再沿电子运动方向求包络的相位，就能把外场作用和能谱边带联系起来。

每一步简化都有自己的条件。正能投影舍去不能真实到达的负能通道；窄带光场选择交换能量；当不同电子通道之间的反冲差别足够小时，最后得到近似均匀的能量梯。下面的关系概括了所用模型之间的联系：
\begin{equation}
{
\text{多模连续谱}
\longrightarrow
\text{正能窄带电子}
\longrightarrow
\text{单频/单模通道}
\longrightarrow
\text{真实反冲格}
\longrightarrow
\text{均匀能量梯}.}
\label{eq:pinem_reduction_spine}
\end{equation}
这些箭头包含对工作能区和场谱的选取，以及对反冲和耦合差别的近似。常见的 Bessel 边带公式将在均匀梯模型中导出；若不同边的失谐或耦合已明显不同，就需要保留各通道的独立演化方程。

先把共振条件写清楚。一次交换若向电子传递机械动量 $\hbar\vb{\kappa}$，电子的能量就从 $E(\vb{p})$ 变为 $E(\vb{p}+\hbar\vb{\kappa})$。设环境同时提供能量 $\hbar\Omega$，两者之差除以 $\hbar$，定义为\emph{失谐}。由于这里采用真实相对论色散，也称为质量壳失谐：
\begin{equation}
{
\Delta_{\rm ms}(\vb{p};\vb{\kappa},\Omega)
\equiv
\frac{
E(\vb{p}+\hbar\vb{\kappa})-E(\vb{p})
}{\hbar}
-\Omega .}
\label{eq:free_electron_exact_mass_shell_mismatch}
\end{equation}
其中 $\vb{\kappa}$ 的单位为 ${\rm m^{-1}}$，$\Omega$ 与 $\Delta_{\rm ms}$ 的单位为 ${\rm s^{-1}}$；$\Omega$ 可取正负，分别表示向电子提供或从电子取走能量。$\Delta_{\rm ms}=0$ 表示能量恰好匹配。根据\chapref{chap:electron-finite-time}的有限时间跃迁公式，作用时间为 $T$ 时，
满足 $|\Delta_{\rm ms}|T\lesssim1$ 的通道位于共振主峰附近。对重复的单一交换步长，
\[
\vb{p}_\ell=\vb{p}_0+\ell\hbar\vb{\kappa},
\qquad
\Delta_\ell
=
\Delta_{\rm ms}(\vb{p}_\ell;\vb{\kappa},\omega),
\]
所以“低反冲”真正要求的是实际占据窗口内 $\Delta_\ell$ 和逐边矩阵元的变化小于动力学分辨尺度；它不等价于只要求 $\hbar\omega/E_0$ 很小。反过来，当不同 $\Delta_\ell$ 可被有限时间分辨时，真实反冲会选择少数共振边，平移对称能量梯必须退出。这一质量壳失谐把 PINEM、有限反冲格和少态共振统一成同一个运动学问题。

\begin{proposition}[孤立自由电子的单光子吸收禁戒]
设电子的入射与出射四动量 $p_i,p_f$ 均位于正能质量壳，
$p_i^2=p_f^2=m_{\mathrm e}^2c^2$；设真实真空光子的四动量 $k$ 满足
$k^2=0$ 且 $\omega>0$。若没有结构、介质或其他粒子参与四动量交换，
则四动量守恒 $p_f=p_i+k$ 无解，即不能发生孤立电子吸收单个真实光子的在壳过程。
\end{proposition}
\begin{proof}
对孤立自由电子吸收一个真实真空光子，若 $p_f=p_i+k$，则 $p_i^2=p_f^2=m_{\mathrm e}^2c^2$ 与 $k^2=0$ 要求 $2p_i\cdot k=0$。但正能电子与正能光子满足
\begin{equation*}
 p_i\cdot k=\frac{E_i\hbar\omega}{c^2}
 (1-\beta_{\rm rel}\cos\vartheta)>0,
\end{equation*}
其中 $\vartheta$ 是入射电子和光子的夹角，$\beta_{\rm rel}=v_i/c<1$。因此无其他参与者的在壳单光子吸收不能发生。近场同步积分没有违背这一点：结构或介质参与了动量守恒。单纯把真空光场人为截成有限空间窗口，不能替代对截断边缘所代表物理源的说明。
\end{proof}

为计算实际电子的速度和反冲，沿用中心动量 $p_0$，相应总能量、Lorentz 因子与群速度为
\begin{equation*}
 p_0,\qquad E_0=\sqrt{m_{\mathrm e}^2c^4+c^2p_0^2},\qquad
 \gamma_0=\frac{E_0}{m_{\mathrm e} c^2},\qquad
 v_0=\frac{c^2p_0}{E_0},
\end{equation*}
纵向二阶色散质量固定写为 $m_\parallel=\gamma_0^3m_{\mathrm e}$。对选中的交换角频率 $\omega$，理想能量梯用 $|\ell\rangle$ 表示，$\ell\in\ZZ$ 是电子相对于中心通道交换的净量子数；真实通道能量始终写 $\mathcal E_\ell=E(p_\ell)$，只有在均匀梯近似成立后才写 $E_\ell^{(0)}=E_0+\ell\hbar\omega$。经典场作用在整个路径上积累的无量纲耦合记为 $\beta$。它将由沿轨迹的场积分确定；$|\beta|$ 控制边带分布的宽度，$\arg\beta$ 控制各边带振幅的相位。光场量子化后的耦合与算符将在量子交换一节引入。

为区别时间域实场与单频复振幅，实时间域电场记为 $\vb{E}(\vb{r},t)$，单频正频分量记为
\begin{equation*}
 \vb{E}^{(+)}(\vb{r},t)=\vb{E}_\omega(\vb{r})\ee^{-\ii\omega t},
 \qquad
 E_{\omega,z}(\vb{r})\equiv\hat{\vb{z}}\cdot\vb{E}_\omega(\vb{r}).
\end{equation*}
$E_{\omega,z}$ 的单位为 $\mathrm{V\,m^{-1}}$。实场由正频分量及其共轭相加，所以纵向实场峰值为 $2|E_{\omega,z}|$。它与单位为能量的通道能量 $\mathcal E_\ell$ 是不同的量。对一般非单色信号，Fourier 变换仍按全书约定用波浪号表示。

离散通道能否近似成均匀能量梯，首先由电子色散曲率决定。若单次交换对应的纵向动量步长为 $\hbar \kappa_z$，二阶色散自然给出反冲频率与有限时间累积相位
\begin{equation}
{
\omega_{\rm r}\equiv\frac{\hbar \kappa_z^2}{2m_\parallel},\qquad
\varphi_{\rm r}\equiv\omega_{\rm r}T,}
\label{eq:pinem_recoil_frequency_master}
\end{equation}
其中 $T$ 是实际作用时间。$\omega_{\rm r}$ 是相邻动量步长经纵向色散曲率产生的反冲频率尺度，$\varphi_{\rm r}$ 是作用窗口内累计的反冲相位。同一 $\omega_{\rm r}$ 与 $\varphi_{\rm r}$ 同时控制有限反冲格、自由漂移相位和实验反演中的反冲分辨，因此这些对象使用同一组定义。

电子--量子场联合哈密顿量的经典场极限由\eqnrefs{eq:multimode_coherent_classical_scaling} 给出：在保持平均电磁场有限时令相对量子涨落趋于零，光场对电子便退化为给定的经典背景。该极限只约束光场的相对量子涨落；电子动量宽度、自旋、反冲、场频谱和作用强度仍由电子侧的独立控制参数决定，因此起点仍是完整外场 Dirac 传播子。

\subsection{外场 Dirac 传播子}

给定外场后，电子波函数满足含电磁耦合的 Dirac 方程。这里把光场视为已知函数，电子的自旋和相对论运动仍由四分量波函数描述。采用 SI 四势
\begin{equation*}
A^\mu(x)\equiv\left(\frac{\phi(\vb{r},t)}{c},\vb{A}(\vb{r},t)\right)
\end{equation*}
参数化外场，其中 $\phi$ 与 $\vb{A}$ 分别是标势和矢势；电子电荷为 $q_{\mathrm e}=-e$。最小耦合 Dirac 哈密顿量由\eqnrefs{eq:feqo-primer-dirac-em} 定义。给定这一哈密顿量，在共同稠密定义域与通常保证非自治演化存在的连续性条件下，定义唯一幺正传播子 $\hat U_A(t_f,t_i)$，满足
\begin{equation}
{
\ii\hbar\partial_{t_f}\hat U_A(t_f,t_i)
=\hat H_D[A](t_f)\hat U_A(t_f,t_i),
\qquad
\hat U_A(t_i,t_i)=\mathbb I.}
\label{eq:exact_external_field_evolution}
\end{equation}
给定入射态后，作用这个传播子就得到出射态；后面的相位调制公式将由它在窄波包条件下的近似求得。

为了区分自由传播与外场作用，写
\begin{equation*}
\hat H_0=c\boldsymbol\alpha\cdot\hat{\vb{p}}+\beta_Dm_{\mathrm e} c^2,
\qquad
\hat V_A(t)=q_{\mathrm e}\phi(\hat{\vb{r}},t)-q_{\mathrm e} c\boldsymbol\alpha\cdot\vb{A}(\hat{\vb{r}},t),
\end{equation*}
并以 $t_i$ 为基准定义相互作用绘景中的作用算符
\begin{equation*}
\hat V_I(t)
\equiv
\ee^{+\ii\hat H_0(t-t_i)/\hbar}
\hat V_A(t)
\ee^{-\ii\hat H_0(t-t_i)/\hbar}.
\end{equation*}
Dyson 时间排序指数已经由 Volterra 积分方程的逐次迭代确定。把当前 Dirac 外场哈密顿量代入这一一般传播子结构，严格得到
\begin{equation}
{
\hat U_A(t_f,t_i)
=\ee^{-\ii\hat H_0(t_f-t_i)/\hbar}\hat U_I(t_f,t_i),
\qquad
\hat U_I
=\mathcal T\exp\!\left[-\frac{\ii}{\hbar}
\int_{t_i}^{t_f}\dd t\,\hat V_I(t)\right].}
\label{eq:exact_external_field_interaction_picture}
\end{equation}
时间排序指数定义完整的有限时传播子；弱作用或短时模型由前文 Dyson/Magnus 展开的截断阶数与相应误差控制量决定。

在坐标表象中定义 Dirac 传播核
\begin{equation}
{
\Psi_f(\vb{r}_f)
=\int\dd^3 r_i\,
K_A(\vb{r}_f,t_f;\vb{r}_i,t_i)\Psi_i(\vb{r}_i),
\qquad
K_A\equiv\langle\vb{r}_f|\hat U_A|\vb{r}_i\rangle .}
\label{eq:exact_external_field_kernel}
\end{equation}
如果入射与出射区域外场已经消失，则可用自由 Dirac 正能态定义正能散射块
\begin{equation}
{
\mathcal S^{++}_{s_fs_i}(\vb{p}_f,\vb{p}_i;t_f,t_i)
=\langle\vb{p}_f,s_f,+|\hat U_A|\vb{p}_i,s_i,+\rangle .}
\label{eq:exact_positive_energy_external_scattering_kernel}
\end{equation}
任意初始动量--自旋振幅因此变成
\begin{equation}
{
a^{(f)}_{s_f}(\vb{p}_f)
=\sum_{s_i}\int\dd^3 p_i\,
\mathcal S^{++}_{s_fs_i}(\vb{p}_f,\vb{p}_i)
a^{(i)}_{s_i}(\vb{p}_i).}
\label{eq:exact_external_field_wavepacket_scattering}
\end{equation}
$\mathcal S^{++}$ 只是完整幺正传播子的一个子块，一般并不单独幺正。对正能初态，概率缺口恰好等于传播到负能子块的概率；若这一泄漏不可忽略，就必须保留完整 $\mathcal S^{\sigma_f\sigma_i}$，而不能把缺口解释成唯象损耗。

\begin{proposition}[幺正传播的正能压缩与泄漏]
设正、负能投影 $P_\pm$ 完备互补，完整外场传播 $U_A$ 幺正。正能入射波包的正能输出由\eqnrefs{eq:exact_external_field_wavepacket_scattering}给出，压缩核满足 $(S^{++})^\dagger S^{++}+(S^{-+})^\dagger S^{-+}=I_+$。因此压缩传播一般仅为收缩，忽略负能支需另控制泄漏。
\end{proposition}
\begin{proof}
把正能入射波包传播后投影到末态，便可同时求出正能振幅与被舍去的负能概率。自由 Dirac 哈密顿量的广义本征态按正、负能两支完备；采用连续归一化，其完备关系为

\begin{equation*}
\mathbb I
=\sum_s\int\dd^3 p\,
\Bigl(
|\vb{p},s,+\rangle\langle\vb{p},s,+|
+|\vb{p},s,-\rangle\langle\vb{p},s,-|
\Bigr)
\equiv P_++P_-.
\end{equation*}
一个纯正能入射波包因此是
\begin{equation*}
|\Psi_i\rangle
=\sum_{s_i}\int\dd^3 p_i\,
 a^{(i)}_{s_i}(\vb{p}_i)|\vb{p}_i,s_i,+\rangle,
\qquad
\sum_{s_i}\int\dd^3 p_i\,
|a^{(i)}_{s_i}(\vb{p}_i)|^2=1.
\end{equation*}
完整外场传播后插入末态完备关系，并投影到正能支：
\begin{align*}
P_+|\Psi_f\rangle
&=P_+U_AP_+|\Psi_i\rangle\\
&=\sum_{s_f}\int\dd^3 p_f\,
|\vb{p}_f,s_f,+\rangle
\langle\vb{p}_f,s_f,+|U_AP_+|\Psi_i\rangle.
\end{align*}
把入射波包展开代回最后一个矩阵元，得到
\begin{align*}
\langle\vb{p}_f,s_f,+|U_A|\Psi_i\rangle
&=\sum_{s_i}\int\dd^3 p_i\,
\langle\vb{p}_f,s_f,+|U_A|\vb{p}_i,s_i,+\rangle
 a^{(i)}_{s_i}(\vb{p}_i)\\
&=\sum_{s_i}\int\dd^3 p_i\,
\mathcal S^{++}_{s_fs_i}(\vb{p}_f,\vb{p}_i)
 a^{(i)}_{s_i}(\vb{p}_i),
\end{align*}
这就是\eqnrefs{eq:exact_external_field_wavepacket_scattering}。积分测度、Dirac $\delta$ 归一化与自旋求和均由自由态基底固定，因此不会额外产生唯象态密度因子。再说明为什么 $S^{++}$ 一般不单独幺正。完整传播子满足 $U_A^\dagger U_A=\mathbb I$，在正能入射子空间两侧夹上 $P_+$：
\begin{align*}
P_+
&=P_+U_A^\dagger(P_++P_-)U_AP_+\\
&=(P_+U_A^\dagger P_+)(P_+U_AP_+)
 +(P_+U_A^\dagger P_-)(P_-U_AP_+).
\end{align*}
于是作为正能子空间上的算符恒等式，
\begin{equation*}
(S^{++})^\dagger S^{++}+(S^{-+})^\dagger S^{-+}=I_+.
\end{equation*}
对任意归一化正能初态取期望值，就得到
\begin{equation*}
\|S^{++}\psi\|^2
=1-\|S^{-+}\psi\|^2.
\end{equation*}
\end{proof}


同一结构同时固定轨迹相位的规范变换性质。采用
\begin{equation}
{
A_\mu' = A_\mu-\partial_\mu\chi,
\qquad
\Psi' = \ee^{\ii q_{\mathrm e}\chi/\hbar}\Psi.}
\label{eq:external_field_gauge_transformation}
\end{equation}
完整传播核满足端点协变律
\begin{equation}
{
K_{A-\partial\chi}(x_f,x_i)
=\ee^{\ii q_{\mathrm e}\chi(x_f)/\hbar}
K_A(x_f,x_i)
\ee^{-\ii q_{\mathrm e}\chi(x_i)/\hbar}.}
\label{eq:exact_dirac_kernel_gauge_covariance}
\end{equation}

\begin{theorem}[时间依赖规范变换的端点协变性]
设实光滑规范函数 $\chi$ 的乘法幺正 $G_\chi$ 保持 Dirac 算符共同域，且两规范下传播子存在。按本章规范约定，$H[A\prime]=G_\chi H[A]G_\chi^\dagger+\ii\hbar\dot G_\chi G_\chi^\dagger$，故传播子和坐标核满足\eqnrefs{eq:exact_dirac_kernel_gauge_covariance}的两端点协变律。
\end{theorem}
\begin{proof}
先在算符层面证明，而不是直接比较两个积分核。外场 Dirac 方程写成
\[
\ii\hbar\partial_t|\Psi(t)\rangle
=
\hat H[A](t)|\Psi(t)\rangle .
\]
规范变换取正文约定
\[
A_\mu'(x)=A_\mu(x)-\partial_\mu\chi(x),
\qquad
|\Psi'(t)\rangle
=
\hat G_\chi(t)|\Psi(t)\rangle,
\]
其中在坐标表象
\[
\langle\vb{r}|\hat G_\chi(t)|\Psi\rangle
=
\ee^{+\ii q_{\mathrm e}\chi(\vb{r},t)/\hbar}
\Psi(\vb{r},t).
\]
对变换后的态求时间导数：
\begin{align*}
\ii\hbar\partial_t|\Psi'\rangle
&=
\ii\hbar(\partial_t\hat G_\chi)|\Psi\rangle
+\hat G_\chi\hat H[A]|\Psi\rangle\\
&=
\left[
\hat G_\chi\hat H[A]\hat G_\chi^\dagger
+\ii\hbar(\partial_t\hat G_\chi)\hat G_\chi^\dagger
\right]
|\Psi'\rangle.
\end{align*}
因此变换后的哈密顿量必须满足
\[
\hat H[A']
=
\hat G_\chi\hat H[A]\hat G_\chi^\dagger
+\ii\hbar(\partial_t\hat G_\chi)\hat G_\chi^\dagger.
\]
第二项正是时间依赖幺正变换中必需的附加项；若只写
\(\hat H'\!=\hat G\hat H\hat G^\dagger\)，会漏掉标势的规范变化。

现在令 \(\hat U_A(t,t_i)\) 是原哈密顿量的时间序传播算符，
\[
\ii\hbar\partial_t\hat U_A(t,t_i)
=
\hat H[A](t)\hat U_A(t,t_i),
\qquad
\hat U_A(t_i,t_i)=I.
\]
定义
\[
\widetilde U(t,t_i)
\equiv
\hat G_\chi(t)\hat U_A(t,t_i)\hat G_\chi^\dagger(t_i).
\]
对 \(t\) 求导，
\begin{align*}
\ii\hbar\partial_t\widetilde U
&=
\left[
\ii\hbar(\partial_t\hat G_\chi)\hat G_\chi^\dagger
+\hat G_\chi\hat H[A]\hat G_\chi^\dagger
\right]\widetilde U\\
&=
\hat H[A'](t)\widetilde U.
\end{align*}
同时
\[
\widetilde U(t_i,t_i)=I.
\]
由传播子初值问题的唯一性，
\[
{
\hat U_{A-\partial\chi}(t_f,t_i)
=
\hat G_\chi(t_f)
\hat U_A(t_f,t_i)
\hat G_\chi^\dagger(t_i).
}
\]

最后取坐标表象。定义
\[
K_A(x_f,x_i)
\equiv
\langle\vb{r}_f|
\hat U_A(t_f,t_i)
|\vb{r}_i\rangle .
\]
利用 \(\hat G_\chi\) 在坐标表象中为乘法算符，
\begin{align*}
K_{A-\partial\chi}(x_f,x_i)
&=
\langle\vb{r}_f|
\hat G_\chi(t_f)\hat U_A
\hat G_\chi^\dagger(t_i)
|\vb{r}_i\rangle\\
&=
\ee^{+\ii q_{\mathrm e}\chi(x_f)/\hbar}
K_A(x_f,x_i)
\ee^{-\ii q_{\mathrm e}\chi(x_i)/\hbar}.
\end{align*}
因此
\[
{
K_{A-\partial\chi}(x_f,x_i)
=
\ee^{\ii q_{\mathrm e}[\chi(x_f)-\chi(x_i)]/\hbar}
K_A(x_f,x_i),
}
\]
即正文\eqnrefs{eq:exact_dirac_kernel_gauge_covariance}。
\end{proof}

\eqnrefs{eq:exact_dirac_kernel_gauge_covariance} 已固定传播核在规范变换下的端点相位。对传播核做正能、慢包络或程函约化时，这个端点变换律必须由约化核继承，因此 Wilson 线与程函相位的规范因子由母传播核的协变性唯一决定。任何降维都必须同时记录被投影的 Hilbert 空间扇区、被舍弃的谱相位以及相应无量纲控制参数。

端点相位律要求任何由同一外场 Dirac 动力学约化得到的传播核继承相同的规范协变结构；Wilson 线与程函近似中的端点相位正是这一要求的具体表现。

正能散射概率的缺口严格等于负能扇区泄漏；只有当 $P_-U_AP_+$ 在实际波包上可忽略时，$S^{++}$ 才能作为近似幺正的有效电子散射算符。

\subsection{正能包络与程函近似}

对于动量集中在 $\vb{p}_0$ 附近的入射波包，大部分空间振荡来自中心平面波。把这一已知振荡因子提出后，剩下的包络变化较慢。若外场主要改变沿途相位，而对前进方向和自旋的影响较小，包络方程可以进一步简化。这种沿近似轨迹积累相位的方法称为\emph{程函近似}（程函近似）。下面从 Dirac 方程检验它的条件，并求出包络满足的方程。

这里的\emph{程函近似}（程函近似）是一种高能、慢变的输运约化：把完整波动方程中快速变化的自由载波显式提出以后，只保留慢包络沿中心群速度轨迹的一阶输运；外场的主要作用因而表现为沿该轨迹积累的相位。它保留相对论色散本身；受控忽略的是自旋子随动量的显著变化、强横向偏转、高阶色散以及正负能扇区混合。相应控制参数先逐一定义，再由传播子误差界判断这一约化在给定电子束和作用长度下是否成立。

这个慢变量就是去掉快速自由载波后的电子包络。设入射波包以 $\vb{p}_0$ 为中心。自由色散关系
\[
E(\vb{p})=\sqrt{m_{\mathrm e}^2c^4+c^2\vb{p}^2}
\]
在 $\vb{p}_0$ 附近的一阶色散梯度给出群速度，因此决定包络中心的平移；二阶及更高阶导数分别产生波包展宽和更高阶畸变。在外场引起的中心动量偏移与自旋旋转均由相应小参数控制时，最低阶输运方程沿群速度轨迹积分，得到经典轨迹上的相位累积。

对窄正能波包，以中心动量 $\vb{p}_0$ 的正能自由旋量作为局域载波基底，并把快速自由相位与慢包络分开：
\begin{equation*}
\Psi(\vb{r},t)
=u_0\,f_{\mathrm e}(\vb{r},t)
\exp\!\left[\frac{\ii}{\hbar}(\vb{p}_0\cdot\vb{r}-E_0t)\right]
+\delta\Psi,
\end{equation*}
其中 $E_0=\sqrt{m_{\mathrm e}^2c^4+c^2p_0^2}$，$u_0$ 是归一化正能 Dirac 自旋子，$f_{\mathrm e}$ 是电子包络，$\delta\Psi$ 收集自旋子动量依赖、色散高阶项和外场诱导的负能分量。自由正能流给出群速度
\begin{equation*}
\vb{v}_0=u_0^\dagger c\boldsymbol\alpha u_0
=\nabla_{\vb{p}}E|_{\vb{p}_0}
=\frac{c^2\vb{p}_0}{E_0}.
\end{equation*}

固定轨迹程函近似只在四类变化都慢时成立：正能自旋子在电子带宽内近似不变、横向偏转小、总动量改变量小、背景和包络都缓于 de Broglie 尺度。令 $\Delta p$ 是初始动量带宽，$\Delta p_{\rm int}$ 是相互作用诱导的最大动量变化，$L_A,L_f$ 分别是外场和包络的最短变化尺度。再令 $u_s(\vb{p})$ 为工作偏振的归一化正能旋量，并定义与任意动量依赖相位无关的秩一投影 $\Pi_s(\vb{p})\equiv u_s(\vb{p})u_s^\dagger(\vb{p})$。于是旋量基底变化可以严格量化为
\begin{align*}
\epsilon_u&\equiv
\sup_{|\vb{p}-\vb{p}_0|\le\Delta p}
\|\Pi_s(\vb{p})-\Pi_s(\vb{p}_0)\|,
&\epsilon_\perp&\equiv\frac{|\Delta\vb{p}_{\perp,{\rm int}}|}{p_0},\\
\epsilon_p&\equiv\frac{\Delta p_{\rm int}}{p_0},
&\epsilon_{\rm WKB}&\equiv\max\!\left(\frac1{\kappa_{\rm e0}L_A},\frac1{\kappa_{\rm e0}L_f}\right).
\end{align*}
其中 $\kappa_{\rm e0}=p_0/\hbar$。方便的局域充分条件是
\begin{equation}
{
\epsilon_{\rm eik}
\equiv\max(\epsilon_u,\epsilon_\perp,\epsilon_p,\epsilon_{\rm WKB})\ll1.}
\label{eq:general_eikonal_control_parameter}
\end{equation}
这还必须配合正负能扇区之间的实际矩阵元远小于约 $2m_{\mathrm e} c^2$ 的谱隙；“电子很快”本身不能替代这一检查。

更可靠的判据是传播子误差。令 $P_+$、$P_-$ 是自由 Dirac 正负能投影，$P_{\mathcal W}\le P_+$ 选取实际电子窄带；以线性化正能包络生成元和自由负能块构造块对角近似哈密顿量$H_{\rm app}$，并定义
\begin{equation*}
R_{\rm eik}(t)=H_D[A](t)-H_{\rm app}(t).
\end{equation*}
则完整 Dirac 传播子 $U_A$ 与近似传播子 $U_{\rm app}$ 满足
\begin{equation}
{
\|[U_A-U_{\rm app}]P_{\mathcal W}\|
\le
\epsilon_{\rm eik}^{\rm op}
\equiv\frac1\hbar\int_{t_i}^{t_f}\dd t\,
\|R_{\rm eik}(t)U_{\rm app}(t,t_i)P_{\mathcal W}\|.}
\label{eq:eikonal_duhamel_error_bound}
\end{equation}
因此真正需要小的是累计余项 $\epsilon_{\rm eik}^{\rm op}$；\eqnrefs{eq:general_eikonal_control_parameter} 只是对它的局域尺度估计。

\begin{proposition}[固定轨迹包络的传播误差]
设完整 Dirac 传播和选定程函参考传播在同一 Hilbert 空间中存在且幺正，余项乘参考演化后的入射窗有定义且范数可积。则\eqnrefs{eq:eikonal_duhamel_error_bound}成立；它同时控制高阶色散、旋量变化、横向偏转及正负能混合的实际累计余项，局部窄带条件须转化为此积分。
\end{proposition}
\begin{proof}
从完整最小耦合 Dirac 方程出发，先把快速自由载波与慢包络分开。取
\[
\Psi(\vb{r},t)
=u_0\,f_{\mathrm e}(\vb{r},t)
\exp\!\left[\frac{\ii}{\hbar}(\vb{p}_0\cdot\vb{r}-E_0t)\right]
+\delta\Psi,
\]
其中自由正能旋量满足
\[
\left(c\boldsymbol\alpha\cdot\vb{p}_0+\beta_Dm_{\mathrm e} c^2\right)u_0=E_0u_0,
\qquad u_0^\dagger u_0=1.
\]
对载波求导时
\begin{align*}
\ii\hbar\partial_t(f_{\mathrm e}\ee^{\ii(\vb{p}_0\cdot\vb{r}-E_0t)/\hbar})
&=\left(E_0f_{\mathrm e}+\ii\hbar\partial_tf_{\mathrm e}\right)
\ee^{\ii(\vb{p}_0\cdot\vb{r}-E_0t)/\hbar},\\
\hat{\vb{p}}(f_{\mathrm e}\ee^{\ii(\vb{p}_0\cdot\vb{r}-E_0t)/\hbar})
&=\left(\vb{p}_0f_{\mathrm e}-\ii\hbar\nabla f_{\mathrm e}\right)
\ee^{\ii(\vb{p}_0\cdot\vb{r}-E_0t)/\hbar}.
\end{align*}
代入
\[
\ii\hbar\partial_t\Psi=
\left[c\boldsymbol\alpha\cdot(\hat{\vb{p}}-q_{\mathrm e}\vb{A})
+\beta_Dm_{\mathrm e} c^2+q_{\mathrm e}\phi\right]\Psi
\]
以后，自由本征方程把所有与 $E_0$ 和 $\vb{p}_0$ 成正比的载波项精确抵消，得到
\[
\ii\hbar u_0\partial_tf_{\mathrm e}
=-\ii\hbar c\boldsymbol\alpha\cdot(\nabla f_{\mathrm e})u_0
+q_{\mathrm e}(\phi-c\boldsymbol\alpha\cdot\vb{A})u_0f_{\mathrm e}
+\mathcal R_D.
\]
$\mathcal R_D$ 收集没有进入固定旋量、固定中心轨迹一阶输运方程的项，包括旋量随动量的变化、横向偏转、高阶色散与正负能混合。左乘 $u_0^\dagger$，并使用自由 Dirac 流恒等式
\[
u_0^\dagger c\boldsymbol\alpha u_0=\vb{v}_0,
\qquad
\vb{v}_0=\nabla_{\vb{p}}E(\vb{p}_0),
\]
可得
\[
\ii\hbar(\partial_t+\vb{v}_0\cdot\nabla)f_{\mathrm e}
=q_{\mathrm e}(\phi-\vb{v}_0\cdot\vb{A})f_{\mathrm e}
+u_0^\dagger\mathcal R_D.
\]
把最后一项删除所定义的近似生成元记为 $H_{\rm app}(t)$，完整 Dirac 生成元与它的差为
\[
R_{\rm eik}(t)=H_D[A](t)-H_{\rm app}(t).
\]
令 $U_A$ 与 $U_{\rm app}$ 分别满足
\begin{align*}
\ii\hbar\partial_tU_A(t,t_i)&=H_D[A](t)U_A(t,t_i),\\
\ii\hbar\partial_tU_{\rm app}(t,t_i)&=H_{\rm app}(t)U_{\rm app}(t,t_i).
\end{align*}
为直接比较两个传播子，定义
\[
F(t)=U_A(t_f,t)U_{\rm app}(t,t_i).
\]
反向传播子的导数满足
\[
\partial_tU_A(t_f,t)=\frac{\ii}{\hbar}U_A(t_f,t)H_D[A](t),
\]
所以乘积求导给
\begin{align*}
\dot F(t)
&=\frac{\ii}{\hbar}U_A(t_f,t)H_D[A](t)U_{\rm app}(t,t_i)
-\frac{\ii}{\hbar}U_A(t_f,t)H_{\rm app}(t)U_{\rm app}(t,t_i)\\
&=\frac{\ii}{\hbar}U_A(t_f,t)R_{\rm eik}(t)U_{\rm app}(t,t_i).
\end{align*}
端点分别为
\[
F(t_f)=U_{\rm app}(t_f,t_i),
\qquad
F(t_i)=U_A(t_f,t_i).
\]
因此
\[
U_A(t_f,t_i)-U_{\rm app}(t_f,t_i)
=-\frac{\ii}{\hbar}\int_{t_i}^{t_f}\dd t\,
U_A(t_f,t)R_{\rm eik}(t)U_{\rm app}(t,t_i).
\]
右乘实际工作窗投影 $P_{\mathcal W}$，取算符范数并使用
$\|U_A(t_f,t)\|=1$ 与范数次乘性，得到
\begin{align*}
\|[U_A-U_{\rm app}]P_{\mathcal W}\|
&\le\frac1\hbar\int_{t_i}^{t_f}\dd t\,
\|U_A(t_f,t)R_{\rm eik}(t)U_{\rm app}(t,t_i)P_{\mathcal W}\|\\
&\le\frac1\hbar\int_{t_i}^{t_f}\dd t\,
\|R_{\rm eik}(t)U_{\rm app}(t,t_i)P_{\mathcal W}\|,
\end{align*}
即正文\eqnrefs{eq:eikonal_duhamel_error_bound}。
\end{proof}


\subsection{轨迹相位与频谱卷积}

在程函近似误差受控以后，正能包络沿经典特征线 $\vb{r}_e(t)=\vb{r}_{e0}+\vb{v}_0t$ 满足一阶标量方程
\begin{equation}
{
\ii\hbar\frac{\dd f_{\mathrm e}}{\dd t}
=q_{\mathrm e}[\phi(\vb{r}_e,t)-\vb{v}_0\cdot\vb{A}(\vb{r}_e,t)]f_{\mathrm e}.}
\label{eq:general_gauge_covariant_eikonal}
\end{equation}
它的解是沿电子世界线的阿贝尔 Wilson 相位
\begin{equation}
{
\frac{f_{\mathrm e,f}}{f_{\mathrm e,i}}
=\exp\!\left[-\frac{\ii q_{\mathrm e}}{\hbar}
\int_{\mathcal C}\dd t\,(\phi-\vb{v}_0\cdot\vb{A})\right].}
\label{eq:general_eikonal_wilson_phase}
\end{equation}
其中 $\mathcal C$ 表示实际电子轨迹。规范变换只改变端点相位；入射、出射基底作相同规范变换后，散射概率保持不变。

令 $\zeta$ 标记电子切片，定义轨迹相位 $\varphi_{\rm eik}(\zeta)$ 和乘法算符 $\mathcal U_{\rm eik}(\zeta)=\ee^{\ii\varphi_{\rm eik}(\zeta)}$。坐标空间的纯相位乘法在动量空间成为卷积：
\begin{equation*}
\widetilde f_{\mathrm e,f}(\kappa)=\int\dd \kappa'\,\mathcal K_{\rm eik}(\kappa-\kappa')\widetilde f_{\mathrm e,i}(\kappa'),
\qquad
\mathcal K_{\rm eik}(\Delta\kappa)=\frac1{2\pi}\int\dd\zeta\,\ee^{-\ii\Delta\kappa\zeta}\mathcal U_{\rm eik}(\zeta).
\end{equation*}
若 $\mathcal U_{\rm eik}$ 对纵向坐标具有周期 $2\pi/\kappa_{\rm g}$，写傅里叶系数 $A_\ell$ 后得到离散平移律
\begin{equation}
{
\widetilde f_{\mathrm e,f}(\kappa)=\sum_{\ell\in\ZZ}A_\ell\widetilde f_{\mathrm e,i}(\kappa-\ell\kappa_{\rm g}).}
\label{eq:general_periodic_eikonal_discrete_convolution}
\end{equation}
周期近场因此不是“产生新的电子能级”，而是其傅里叶谐波选择了一组相差 $\hbar\kappa_{\rm g}$ 的自由电子动量副本。

\begin{proposition}[实周期程函相位的动量卷积]
沿给定特征线，实标量输运方程的解是\eqnrefs{eq:general_eikonal_wilson_phase}的单位模相位。若该相位因子以 $2\pi/\kappa_{\rm g}$ 为周期，其 Fourier 系数给\eqnrefs{eq:general_periodic_eikonal_discrete_convolution}的离散动量卷积，系数平方和为一。势改写成电场积分时须保留或明确消去端点项。
\end{proposition}
\begin{proof}
其一般系数给出正文\eqnrefs{eq:general_periodic_eikonal_discrete_convolution}，而单色正弦相位的 Bessel 系数只是这一 Fourier 系数的特例。 沿固定特征线 $\vb{r}_e(t)=\vb{r}_{e0}+\vb{v}_0t$，链式法则给

\begin{equation*}
\frac{\dd}{\dd t}f_{\mathrm e}[\vb{r}_e(t),t]
=(\partial_t+\vb{v}_0\cdot\nabla)f_{\mathrm e}.
\end{equation*}
因此正能包络方程化为一阶常微分方程
\begin{equation*}
\frac{1}{f_{\mathrm e}}\frac{\dd f_{\mathrm e}}{\dd t}
=-\frac{\ii q_{\mathrm e}}{\hbar}
[\phi(\vb{r}_e,t)-\vb{v}_0\cdot\vb{A}(\vb{r}_e,t)].
\end{equation*}
把两边对时间积分，左边是对数的全微分：
\begin{align*}
\int_{t_i}^{t_f}\dd t\,\frac1f_{\mathrm e}\frac{\dd f_{\mathrm e}}{\dd t}
&=\ln\frac{f_{\mathrm e}(t_f)}{f_{\mathrm e}(t_i)},\\
&=-\frac{\ii q_{\mathrm e}}{\hbar}\int_{t_i}^{t_f}\dd t\,
[\phi(\vb{r}_e,t)-\vb{v}_0\cdot\vb{A}(\vb{r}_e,t)].
\end{align*}
指数化后得到沿电子世界线的 Wilson 相位因子
\begin{equation*}
\frac{f_{\mathrm e}(t_f)}{f_{\mathrm e}(t_i)}
=\exp\!\left\{-\frac{\ii q_{\mathrm e}}{\hbar}\int_{t_i}^{t_f}\dd t\,
[\phi-\vb{v}_0\cdot\vb{A}]\right\}.
\end{equation*}
若作规范变换
$A_\mu\to A_\mu-\partial_\mu\chi$，指数中新增
\begin{align*}
-\frac{\ii q_{\mathrm e}}{\hbar}
\int_{t_i}^{t_f}\dd t\,
[-\partial_t\chi-\vb{v}_0\cdot\nabla\chi]
&=\frac{\ii q_{\mathrm e}}{\hbar}
\int_{t_i}^{t_f}\dd t\,\frac{\dd}{\dd t}\chi[\vb{r}_e(t),t]\\
&=\frac{\ii q_{\mathrm e}}{\hbar}[\chi(x_f)-\chi(x_i)],
\end{align*}
正好重现完整 Dirac 传播核的端点协变因子。对单色正频分量写
\begin{equation*}
\phi^{(+)}=\phi_\omega(\vb{r})\ee^{-\ii\omega t},
\qquad
\vb{A}^{(+)}=\vb{A}_\omega(\vb{r})\ee^{-\ii\omega t}.
\end{equation*}
沿直线轨迹取 $z=v_0t$。利用
\begin{equation*}
E_{\omega,z}=-\partial_z\phi_\omega+\ii\omega A_{\omega z}
\end{equation*}
以及
\begin{equation*}
\frac{\dd}{\dd z}
\left[\phi_\omega(z)\ee^{-\ii\omega z/v_0}\right]
=(\partial_z\phi_\omega)\ee^{-\ii\omega z/v_0}
-\frac{\ii\omega}{v_0}\phi_\omega\ee^{-\ii\omega z/v_0},
\end{equation*}
可把势组合整理成
\begin{align*}
\int\dd z\,
\left(\frac{\phi_\omega}{v_0}-A_{\omega z}\right)
\ee^{-\ii\omega z/v_0}
&=\frac{1}{\ii\omega}
\int\dd z\,E_{\omega,z}(z)
\ee^{-\ii\omega z/v_0}\\
&\quad
-\frac{1}{\ii\omega}
\left[\phi_\omega(z)\ee^{-\ii\omega z/v_0}\right]_{-\infty}^{+\infty}.
\end{align*}
对散射问题，作用区外场消失，端点项为零；这就说明同步纵向复幅 $E_{\omega,z}$ 为何给出规范无关的实际耦合来源。最后推导周期相位的动量卷积。若 $\mathcal U_{\rm eik}(\zeta)=\ee^{\ii\varphi_{\rm eik}(\zeta)}$，则
\begin{align*}
\widetilde f_{\mathrm e,f}(\kappa)
&=\int\dd\zeta\,\ee^{-\ii \kappa\zeta}
\mathcal U_{\rm eik}(\zeta)f_{\mathrm e,i}(\zeta)\\
&=\int\frac{\dd \kappa'}{2\pi}
\left[\int\dd\zeta\,
 \ee^{-\ii(\kappa-\kappa')\zeta}\mathcal U_{\rm eik}(\zeta)\right]
\widetilde f_{\mathrm e,i}(\kappa').
\end{align*}
若 $\mathcal U_{\rm eik}$ 的周期为 $2\pi/\kappa_{\rm g}$，展开
\begin{equation*}
\mathcal U_{\rm eik}(\zeta)=\sum_{\ell\in\ZZ}A_\ell \ee^{\ii\ell\kappa_{\rm g}\zeta},
\qquad
A_\ell=\frac{\kappa_{\rm g}}{2\pi}
\int_0^{2\pi/\kappa_{\rm g}}\dd\zeta\,
\mathcal U_{\rm eik}(\zeta)\ee^{-\ii\ell\kappa_{\rm g}\zeta}.
\end{equation*}
代回 Fourier 积分后出现
$2\pi\delta(\kappa-\kappa'-\ell\kappa_{\rm g})$，执行 $\kappa'$ 积分便得到\eqnrefs{eq:general_periodic_eikonal_discrete_convolution}。由于 $|\mathcal U_{\rm eik}|=1$，Parseval 还给
\begin{equation*}
\sum_\ell|A_\ell|^2=1,
\end{equation*}
\end{proof}


要把这些副本解释成互斥边带概率，还必须检查它们近似正交。对归一化入射谱定义
\begin{equation}
{
\epsilon_{\rm ov}
\equiv\sup_{\ell\neq m}
\frac{\left|\int\frac{\dd \kappa}{2\pi}\,
\widetilde f_{\mathrm e,i}^*(\kappa-\ell\kappa_{\rm g})\widetilde f_{\mathrm e,i}(\kappa-m\kappa_{\rm g})\right|}
{\int\frac{\dd \kappa}{2\pi}|\widetilde f_{\mathrm e,i}(\kappa)|^2}\ll1.}
\label{eq:translated_packet_overlap_condition}
\end{equation}
只有在这一条件成立时，$|A_\ell|^2$ 才可直接解释为分辨的边带权重；否则不同平移副本之间的交叉项就是可观测相干。

周期纯相位散射还具有一个比 Bessel 归一化更一般的幺正约束。把
\[
\mathcal U_{\rm eik}(\zeta)
=\sum_{n\in\ZZ}A_n\ee^{+\ii n\kappa_{\rm g}\zeta}
\]
看作单位圆上的 Fourier 级数，并在理想离散通道空间 $\ell^2(\ZZ)$ 上定义 Toeplitz 散射矩阵
\begin{equation}
{
S_{\ell m}=A_{\ell-m}.}
\label{eq:periodic_phase_toeplitz_matrix}
\end{equation}
由于纯相位因子逐点满足 $|\mathcal U_{\rm eik}(\zeta)|=1$，其 Fourier 系数并非任意，而必须满足
\begin{equation}
{
\sum_{n\in\ZZ}A_{n+r}A_n^*=\delta_{r0},
\qquad r\in\ZZ.}
\label{eq:periodic_phase_fourier_unitarity_identity}
\end{equation}
这正是 $S^\dagger S=SS^\dagger=I$ 的矩阵元条件。特别地，$r=0$ 给出
$\sum_n|A_n|^2=1$；$r\neq0$ 则给出不同列之间的正交性。因此任意周期纯相位调制都定义一个严格幺正的双边通道散射，而 Bessel 恒等式
$\sum_nJ_n^2(x)=1$ 只是单谐波相位的特例。若实际平移波包不满足\eqnrefs{eq:translated_packet_overlap_condition}，不能把 $|A_n|^2$ 当成互斥概率，但完整连续波函数的范数仍由这个母幺正性严格保持。

\begin{theorem}[单位模函数诱导双边 Toeplitz 幺正]
设 $U(\theta)$ 为可测 $2\pi$ 周期函数，$|U|=1$ 几乎处处，Fourier 系数为 $A_n$。则 $\sum_m A_{m+r}A_m^*=\delta_{r0}$，每个和绝对收敛；矩阵 $S_{\ell m}=A_{\ell-m}$ 定义 $\ell^2(\ZZ)$ 上的幺正算符。一般输入的卷积按 $\ell^2$ 极限解释，有限边带截断不自动保持严格幺正。
\end{theorem}
\begin{proof}
令无量纲周期相位变量为 $\theta=\kappa_{\rm g}\zeta$，并把一个周期上的纯相位函数写成
\[
U(\theta)=\ee^{\ii\varphi(\theta)}
=\sum_{n\in\ZZ}A_n\ee^{+\ii n\theta},
\qquad
A_n=\frac1{2\pi}\int_{-\pi}^{\pi}\dd\theta\,
U(\theta)\ee^{-\ii n\theta}.
\]
因为 $\varphi(\theta)$ 为实函数，逐点有
\[
U(\theta)U^*(\theta)=1.
\]
把两边都展开成 Fourier 级数：
\begin{align*}
1
&=\left(\sum_nA_n\ee^{+\ii n\theta}\right)
  \left(\sum_mA_m^*\ee^{-\ii m\theta}\right)\\
&=\sum_{n,m}A_nA_m^*\ee^{+\ii(n-m)\theta}.
\end{align*}
令 $r=n-m$，即 $n=m+r$，则
\[
1
=\sum_{r\in\ZZ}
\left(\sum_mA_{m+r}A_m^*\right)
\ee^{+\ii r\theta}.
\]
常数函数 $1$ 的 Fourier 系数只有零阶非零，因此逐阶比较得到
\[
\sum_mA_{m+r}A_m^*=\delta_{r0},
\]
即正文\eqnrefs{eq:periodic_phase_fourier_unitarity_identity}。

现在验证它与离散通道散射矩阵的幺正性完全等价。由
$S_{\ell m}=A_{\ell-m}$，有
\begin{align*}
(S^\dagger S)_{mn}
&=\sum_\ell S_{\ell m}^*S_{\ell n}\\
&=\sum_\ell A_{\ell-m}^*A_{\ell-n}.
\end{align*}
令 $q=\ell-m$，则 $\ell-n=q+m-n$，于是
\[
(S^\dagger S)_{mn}
=\sum_q A_q^*A_{q+m-n}
=\delta_{mn}.
\]
同理 $SS^\dagger=I$。因此任意实周期相位都生成一个双边 Toeplitz 幺正算符；这一结论与相位是否正弦、驱动是否弱、边带数是否有限都无关。

若入射态在离散通道基中为 $|\psi_i\rangle=\sum_m c_m|m\rangle$，输出为
\[
c_\ell'=
\sum_mA_{\ell-m}c_m.
\]
幺正性立即给
\[
\sum_\ell|c_\ell'|^2
=\sum_m|c_m|^2.
\]
对于单通道入射 $c_m=\delta_{m0}$，进一步得到 $P_\ell=|A_\ell|^2$。
\end{proof}


对单色正频场，轨迹相位还可以改写为电场的同步 Fourier 投影。写成
$\vb{E}^{(+)}(\vb{r},t)=\vb{E}_\omega(\vb{r})\ee^{-\ii\omega t}$，并定义
$\vb{E}_\omega=-\nabla\phi_\omega+\ii\omega\vb{A}_\omega$。这样，轨迹上的势组合可以分成纵向电场积分与端点势差。散射边界条件要求电子进入和离开作用区时外场趋于零，端点项随之消失。因此 PINEM 累积耦合统一由沿电子轨迹取样的纵向单频复幅 $E_{\omega,z}$ 表示，其空间依赖同时给出后续所需的场包络。

实际连续电子波包若不同平移副本重叠，则离散通道基不再严格正交，应回到正文\eqnrefs{eq:general_periodic_eikonal_discrete_convolution} 在振幅层保留交叉项；连续 Hilbert 空间中的母传播子仍保持幺正。

在平移副本互相正交时，Parseval 等式说明这些 Fourier 权重组成归一化的边带概率分布。

\subsection{低反冲平移极限}

现在让沿轨迹的场主要以角频率 $\omega$ 振荡。相位调制随电子到达时刻周期变化，它的 Fourier 系数就是能量边带振幅。为了使用同一个系数描述每次交换，需要各条边的耦合近似相同，且反冲造成的相对相位在作用时间内足够小。保留光脉冲的有限包络，把正频纵向场写成
\begin{equation*}
\hat{\vb{z}}\cdot\vb{E}^{(+)}(z,t)
=E_{\omega,z}(z)f(t-\tau)\ee^{-\ii\omega t},
\end{equation*}
其中 $E_{\omega,z}(z)$ 是载频 $\omega$ 的纵向复场幅，$f$ 是无量纲慢变包络，$\tau$ 是场脉冲相对于电子的时间延迟。沿中心速度 $v_0$ 的电子轨迹积累的无量纲单频耦合定义为
\begin{equation}
{
\beta(\zeta,\tau)
=\frac{q_{\mathrm e}}{\hbar\omega}
\int_{-\infty}^{\infty}\dd z\,
E_{\omega,z}(z)
 f\!\left(\frac{z-\zeta}{v_0}-\tau\right)
\ee^{-\ii\omega z/v_0}.}
\label{eq:pinem_standard_beta_pulsed}
\end{equation}
$\zeta$ 标记电子包络内部的纵向切片。$|\beta|$ 衡量沿轨迹相干累积的交换强度，$\arg\beta$ 则保存该切片相对于光学载波的相位。这个积分只使用单频纵向场与给定电子轨迹；代入低反冲平移模型后便得到 PINEM 参数化。

\begin{exbox}[PINEM 低反冲平移极限]
取一维准直电子与单频纵向外场，求从轨迹包络退化到单一 Bessel 边带分布还需要哪些传播条件。

\textbf{解。}
对准直电子，若采用辐射规范并在一维纵向几何中把 $A^2$、横向偏转作为高阶项，常用充分条件是
\begin{equation}
{
\frac{|q_{\mathrm e} A_z|}{p_0}\ll1,
\qquad
\frac{\Delta p_\perp}{p_0}\ll1.}
\label{eq:pinem_ap_small_parameters}
\end{equation}
但高速电子的主方程仍来自 Dirac 正能投影，而不是非相对论 Schr\"odinger 方程。对 $\phi=0$、$\vb{v}_0=v_0\hat z$，
\begin{equation}
{
\ii\hbar(\partial_t+v_0\partial_z)f_{\mathrm e}
\simeq-q_{\mathrm e} v_0A_z(z,t)f_{\mathrm e}.}
\label{eq:dirac_projected_pinem_envelope}
\end{equation}

这里舍去的色散和投影余项具有能量单位，其累计传播效应由\eqnrefs{eq:eikonal_duhamel_error_bound}控制。\eqnrefs{eq:pinem_standard_beta_pulsed}已把有限脉冲与纵向场几何压缩为切片依赖的复耦合 $\beta(\zeta,\tau)$。要进一步得到单一 Bessel 边带序列，还需要电子包络内部不同切片看到近似相同的 $\beta$。
令入射电子纵向包络的显著支持宽度为 $\Delta\zeta_e$。对时间包络 $f(t)$ 定义变化时间
\begin{equation*}
 \tau_f\equiv\frac{\|f\|_\infty}{\|\dot f\|_\infty},
 \qquad
 \epsilon_{\rm env}\equiv\frac{\Delta\zeta_e}{v_0\tau_f}\ll1,
\end{equation*}
其中 $\|\dot f\|_\infty=0$ 时取 $\tau_f=\infty$，全零脉冲则直接取 $\beta=0$。由均值定理与\eqnrefs{eq:pinem_standard_beta_pulsed}，对中心切片 $\zeta_0$ 有绝对误差界
\[
 |\beta(\zeta,\tau)-\beta(\zeta_0,\tau)|
 \le B_{\rm abs}\epsilon_{\rm env},\qquad
 B_{\rm abs}=\frac{e\|f\|_\infty}{\hbar\omega}
                 \int|E_{\omega,z}(z)|\dd z,
\]
这里 $|\zeta-\zeta_0|\le\Delta\zeta_e$，并假定场幅的绝对积分有限。实际同步积分可以因相消而为零，所以这个绝对界不能无条件改写成相对于 $\beta(\zeta_0,\tau)$ 的小误差。局域相位之差至多为 $2B_{\rm abs}\epsilon_{\rm env}$；当这个累计量很小时，可把切片耦合替换为共同值 $\beta=\beta(\zeta_0,\tau)$，输出近似为周期相位
\begin{equation*}
\frac{f_{\mathrm e}^{(f)}(\zeta)}{f_{\mathrm e}^{(i)}(\zeta)}
=\exp[\beta \ee^{\ii\omega\zeta/v_0}-\beta^*\ee^{-\ii\omega\zeta/v_0}].
\end{equation*}
由 Jacobi--Anger 展开，分辨边带的概率为
\begin{equation*}
P_\ell=J_\ell^2(2|\beta|).
\end{equation*}
这一 Bessel 分布同时要求切片耦合的累计相位差、平移对称传播误差和格边界传播误差均足够小；任一条件失效时都应恢复位置依赖耦合或有限反冲动力学。
\end{exbox}
% -----------------------------------------------------------------------------
% -----------------------------------------------------------------------------

\begin{exbox}[200\,keV 电子穿过 800\,nm 高斯近场]
\eqnrefs{eq:pinem_standard_beta_pulsed} 已经把经典 PINEM 的全部场依赖压缩成沿电子轨迹的同步 Fourier 投影。这个积分衡量的是电子轨迹对同步 Fourier 分量的选择。取一个最简单但非平凡的纵向近场包络
\begin{equation*}
 E_{\omega,z}(z)=E_{\omega,0}\exp\!\left(-\frac{z^2}{2L^2}\right),
 \qquad f(t)=1,
\end{equation*}
其中 $E_{\omega,0}$ 是正频电场复振幅的峰值（实场余弦峰值为其模长的两倍），$L$ 是近场沿轨迹的均方根空间宽度。电子初始动能取
\begin{equation*}
 K_0=200\,\mathrm{keV},
\end{equation*}
光学波长、近场峰值和空间宽度分别取
\begin{equation*}
 \lambda=800\,\mathrm{nm},\qquad
 E_{\omega,0}=20\,\mathrm{MV\,m^{-1}},\qquad
 L=100\,\mathrm{nm}.
\end{equation*}

求轨迹积分、完整边带分布和低反冲控制参数，判断局域弱场是否意味着只能保留一阶边带。

\textbf{解。} 先由自由电子相对论运动学得到
\begin{align*}
 \gamma_0
 &=1+\frac{K_0}{m_{\mathrm e} c^2}
 =1.39139,\\
 \frac{v_0}{c}
 &=\sqrt{1-\gamma_0^{-2}}
 =0.69531,\\
 v_0&=2.0845\times10^8\,\mathrm{m\,s^{-1}}.
\end{align*}
光学角频率和单光子能量则为
\begin{align*}
 \omega&=\frac{2\pi c}{\lambda}
 =2.3546\times10^{15}\,\mathrm{s^{-1}},\\
 \hbar\omega&=1.5498\,\mathrm{eV}.
\end{align*}
在这个例子中，\eqnrefs{eq:pinem_standard_beta_pulsed} 变成一个标准高斯 Fourier 积分：
\begin{align}
 \beta
 &=\frac{q_{\mathrm e} E_{\omega,0}}{\hbar\omega}
 \int_{-\infty}^{\infty}\dd z\,
 \exp\!\left(-\frac{z^2}{2L^2}\right)
 \exp\!\left(-\ii\frac{\omega z}{v_0}\right)\notag\\
 &=\frac{q_{\mathrm e} E_{\omega,0}\sqrt{2\pi}L}{\hbar\omega}
 \exp\!\left[-\frac12\left(\frac{\omega L}{v_0}\right)^2\right].
 \label{eq:pinem-gaussian-beta-dp65}
\end{align}
这里
\begin{equation*}
 \varphi_{\rm pm}\equiv\frac{\omega L}{v_0}=1.1296
\end{equation*}
是电子飞越近场包络时累计的光学相位尺度，并由\eqnrefs{eq:pinem_standard_beta_pulsed} 中 Fourier 相位的无量纲组合直接确定。数值代入后
\begin{equation}
 {|\beta|=1.7092,\qquad 2|\beta|=3.4183.}
 \label{eq:pinem-gaussian-beta-number-dp65}
\end{equation}

于是理想低反冲边带概率不需要再解微分方程，直接是
\begin{equation*}
 P_\ell=J_\ell^2(3.4183).
\end{equation*}
几个主要边带为
\begin{center}
\begin{tabular}{c|cccccc}
 $|\ell|$ & 0 & 1 & 2 & 3 & 4 & 5\\\hline
 $P_\ell$ & 0.1351 & 0.02944 & 0.2189 & 0.1413 & 0.03685 & 0.00538
\end{tabular}
\end{center}
这里表中 $|\ell|\ge1$ 的数值是\emph{单侧}边带概率；因为当前纯相位调制满足 $P_{+\ell}=P_{-\ell}$，总概率仍按两侧共同归一化。输出能量峰位为
\begin{equation*}
 K_\ell\simeq K_0+\ell\hbar\omega
 =200\,\mathrm{keV}+1.5498\ell\,\mathrm{eV}.
\end{equation*}
Bessel 权重并不随 $|\ell|$ 单调衰减；例如这里 $|\ell|=2$ 比 $|\ell|=1$ 更强。这不是异常，而是不同量子交换路径在同一整数能量阶上的相干衍射结果。

利用恒等式
\begin{equation*}
 \sum_{\ell=-\infty}^{\infty}\ell^2J_\ell^2(x)=\frac{x^2}{2},
\end{equation*}
还可以不逐项求和就得到电子边带阶数的均方根宽度
\begin{align*}
 \sqrt{\langle\ell^2\rangle}
 &=\sqrt{2}\,|\beta|=2.417,\\
 \sigma_E
 &=\hbar\omega\sqrt{\langle\ell^2\rangle}
 =3.75\,\mathrm{eV}.
\end{align*}
这给出了一个很有用的尺度判断：$|\beta|$ 不是“平均吸收光子数”，而是控制相干能量衍射宽度的脉冲面积。

\begin{center}
\includegraphics[width=0.78\textwidth]{pinem_gaussian_example_dp65.png}\\[-0.2em]
{\small 计算图：200\,keV、800\,nm 高斯近场的离散边带概率。每一根竖线相隔 $\hbar\omega=1.5498\,\mathrm{eV}$，全部权重由同一个 $|\beta|=1.7092$ 决定。}
\end{center}

最后必须检查这组参数是否真的允许使用 Bessel 极限。局域最小耦合小参数为
\begin{equation*}
 \frac{|q_{\mathrm e}|E_{\omega,0}/\omega}{p_0}
 =5.15\times10^{-6},
 \qquad p_0=\gamma_0m_{\mathrm e} v_0,
\end{equation*}
说明单点场强对载波动量只是微扰；但沿整个作用区相干积分后 $|\beta|$ 已经达到 $O(1)$。这两个量衡量的是不同问题，不能用“局域相互作用很弱”推断“最终只会出现一阶边带”。

再检查反冲。同步纵向波数为
\begin{equation*}
 \kappa_z=\frac{\omega}{v_0}=1.1296\times10^7\,\mathrm{m^{-1}},
\end{equation*}
纵向相对论有效质量为 $m_\parallel=\gamma_0^3m_{\mathrm e}=2.6937m_{\mathrm e}$，所以
\begin{align*}
 \omega_{\rm r}
 &=\frac{\hbar\kappa_z^2}{2m_\parallel}
 =2.742\times10^9\,\mathrm{s^{-1}},\\
 \hbar\omega_{\rm r}
 &=1.80\,\mu\mathrm{eV}.
\end{align*}
若用 $T_{\rm int}=2L/v_0=0.959\,\mathrm{fs}$ 作为作用时间尺度，并取实际显著占据窗口 $|\ell|\le5$，则\eqnrefs{eq:pinem_recoil_phase_condition} 给
\begin{equation}
 {
 \epsilon_{\rm rec}^{(T)}
 =(2\times5+1)\omega_{\rm r}T_{\rm int}
 =2.89\times10^{-5}.}
 \label{eq:pinem-gaussian-recoil-number-dp65}
\end{equation}
因此这组参数处在非常干净的低反冲区。反过来，在相同电子能量和光学频率下，如果周期结构把相互作用延长到 $T=1\,\mathrm{ps}$，同一参数已增至约 $3.0\times10^{-2}$；到 $10\,\mathrm{ps}$ 则约为 $0.30$。所以反冲是否重要取决于\emph{累计相位可分辨性}，而不只取决于单次交换的 $\hbar\omega/E_0$。
\end{exbox}


真实自由电子只有正能谱下界，所以理想双边格 $\ell\in\ZZ$ 只是局部坐标近似。若 $\Pi_{\rm edge}$ 投影到最低若干可访问格点，定义
\begin{equation}
{
P_{\rm edge}
=\sup_t\operatorname{Tr}[\rho_e(t)\Pi_{\rm edge}]\ll1.}
\label{eq:bilateral_ladder_edge_condition}
\end{equation}
若累计交换强度为 $\Lambda_I=\int\dd t\,|\upsilon_I(t)|$，则把真实半无限平移算符延拓成双边幺正平移算符的传播误差满足
\begin{equation}
{
\|[U_b-U_s]|\psi_i\rangle\|
\le\Lambda_I\sqrt{P_{\rm edge}},
\qquad
\epsilon_{\rm edge}^{\rm dyn}=\Lambda_I\sqrt{P_{\rm edge}}\ll1.}
\label{eq:bilateral_ladder_dynamical_error}
\end{equation}

在边界和逐边反冲都不可分辨时，引入双边平移算符 $\hat b|\ell\rangle=|\ell-1\rangle$。由于 $[\hat b,\hat b^\dagger]=0$，任意时间依赖的平移对称生成元在不同时刻也对易，因此传播子只依赖复脉冲面积
\begin{equation}
{
S_{\rm tr}(\mathcal Z)
=\exp(\mathcal Z \hat b^\dagger-\mathcal Z^*\hat b).}
\label{eq:universal_translation_propagator}
\end{equation}
在统一的相位基底
$|\theta\rangle=(2\pi)^{-1/2}\sum_\ell \ee^{\ii\ell\theta}|\ell\rangle$
中，
$\hat b|\theta\rangle=\ee^{\ii\theta}|\theta\rangle$、
$\hat b^\dagger|\theta\rangle=\ee^{-\ii\theta}|\theta\rangle$，所以
\begin{equation}
{
S_{\rm tr}(\mathcal Z)|\theta\rangle
=\exp[\mathcal Z \ee^{-\ii\theta}-\mathcal Z^*\ee^{\ii\theta}]|\theta\rangle.}
\label{eq:universal_translation_phase_solution}
\end{equation}
Bessel 边带只是这一相位函数的傅里叶系数。

\begin{proposition}[可交换双边平移生成元的精确传播]
设 $b$ 为 $\ell^2(\ZZ)$ 上的双边单位平移，$\upsilon$ 可积。对 $H_I=\ii\hbar(\upsilon b^\dagger-\upsilon^*b)$，各时刻生成元对易，传播子为 $\exp(\mathcal Zb^\dagger-\mathcal Z^*b)$，$\mathcal Z=\int\upsilon$；其 Fourier 相位为\eqnrefs{eq:universal_translation_phase_solution}。
\end{proposition}
\begin{proof}
正文已经把真实正能谱的下边界误差压缩为\eqnrefs{eq:bilateral_ladder_dynamical_error}。在该误差可忽略、且逐边耦合在实际占据窗口内近似相同时，主线模型只剩双边平移对称生成元
\begin{equation*}
H_I(t)=\ii\hbar[\upsilon(t)\hat b^\dagger-\upsilon^*(t)\hat b].
\end{equation*}
因为 $[\hat b,\hat b^\dagger]=0$，任意两个时刻的生成元满足
\begin{equation*}
[H_I(t),H_I(t')]=0.
\end{equation*}
因此时间排序指数化为普通指数。令
\begin{equation*}
\mathcal Z=\int_{t_i}^{t_f}\dd t\,\upsilon(t),
\end{equation*}
则
\begin{align*}
U_I(t_f,t_i)
&=\exp\!\left[-\frac{\ii}{\hbar}\int_{t_i}^{t_f}\dd t\,H_I(t)\right]\\
&=\exp\!\left[\int_{t_i}^{t_f}\dd t\,
(\upsilon(t)\hat b^\dagger-\upsilon^*(t)\hat b)\right]\\
&=\exp\!\left(\mathcal Z\hat b^\dagger-\mathcal Z^*\hat b\right),
\end{align*}
即正文\eqnrefs{eq:universal_translation_propagator}。

再使用正文给出的 Fourier 相位基
\begin{equation*}
|\theta\rangle=(2\pi)^{-1/2}\sum_{\ell\in\ZZ}
 \ee^{\ii\ell\theta}|\ell\rangle,
\end{equation*}
有
\begin{equation*}
\hat b|\theta\rangle=\ee^{\ii\theta}|\theta\rangle,
\qquad
\hat b^\dagger|\theta\rangle=\ee^{-\ii\theta}|\theta\rangle.
\end{equation*}
把这两个本征值代入指数，得到
\begin{equation*}
S_{\rm tr}(\mathcal Z)|\theta\rangle
=\exp[\mathcal Z \ee^{-\ii\theta}-\mathcal Z^*\ee^{\ii\theta}]|\theta\rangle,
\end{equation*}
即正文\eqnrefs{eq:universal_translation_phase_solution}。
\end{proof}


具体 PINEM 近场确定 $\upsilon(t)$ 或 $\mathcal Z$ 的系统参数，而平移格上的闭式传播子由一般代数结构固定。

\subsection{把场强、相位匹配和边带概率连成一次计算}
\label{sec:electron-pinem-field-calibration}

前面已把电子的调制写成沿轨迹的场积分 $\beta$。现在用一个给定近场求可测的边带概率。代入场幅之前，先把实场峰值与复振幅的关系写清楚：本书的 $E_{\omega,z}$ 表示正频分量，因此单频实场为
\begin{equation}
 E_z(z,t)=E_{\omega,z}(z)\ee^{-\ii\omega t}
           +E_{\omega,z}^*(z)\ee^{\ii\omega t}
 =2|E_{\omega,z}(z)|\cos[\omega t-\arg E_{\omega,z}(z)].
 \label{eq:electron-real-field-positive-frequency}
\end{equation}
因此实场的余弦峰值是 $E_{\rm pk}^{\rm(real)}=2|E_{\omega,z}|$，均方根值是 $\sqrt2|E_{\omega,z}|$。若仿真采用 $E_z=\Re[\mathcal E_{\rm ph}\ee^{-\ii\omega t}]$，应输入 $E_{\omega,z}=\mathcal E_{\rm ph}/2$。漏掉这个二分之一会使 $|\beta|$ 加倍，并使弱边带概率放大四倍。

从\eqnrefs{eq:pinem_standard_beta_pulsed} 出发，先令脉冲包络在电子通过期间近似恒定，再考虑轨迹上长度 $L_{\rm int}$ 内的一个纵向空间谐波：
\begin{equation*}
 E_{\omega,z}(z)=E_{\omega,0}\ee^{\ii\kappa_m z},\quad
 -L_{\rm int}/2<z<L_{\rm int}/2,
 \qquad \Delta\kappa=\kappa_m-\frac{\omega}{v_0}.
\end{equation*}
$E_{\omega,0}$ 是复正频振幅，$\kappa_m$ 是场在轨迹方向的波数；矩形包络是实际场采样的简化模型，其边缘需要由结构或源支持，不能解释为无界真空中的自由平面波。有限长度同步积分给出
\begin{equation}
 \beta=\frac{q_{\mathrm e} E_{\omega,0}L_{\rm int}}{\hbar\omega}
 \sinc\!\left(\frac{\Delta\kappa L_{\rm int}}2\right),\qquad
 P_\ell=J_\ell^2(2|\beta|).
 \label{eq:electron-rectangular-nearfield-beta}
\end{equation}
每一小段 $\dd z$ 都贡献一个复振幅；失配使这些复数在复平面转动，相加后部分抵消。因此更长的作用区只有在相位匹配时才持续增强耦合。取复数而不只取场强平方，是保留这种相消所必需的。

对 $v_0<c$ 的电子，同步波数 $\omega/v_0$ 大于真空光波数 $\omega/c$。均匀真空中的传播光无法提供这个纵向 Fourier 分量；结构近场、介质或倒格矢可以承担所需动量。对于真空侧沿表面传播、垂直表面衰减的单个谐波，Maxwell 色散给
\begin{equation*}
 \kappa_\perp^{\rm(ev)}=\sqrt{\kappa_m^2-\omega^2/c^2},\qquad
 \ell_{\rm ev}=1/\kappa_\perp^{\rm(ev)}.
\end{equation*}
实正的 $\kappa_\perp^{\rm(ev)}$ 表示\emph{倏逝}：场幅随离开表面的距离按指数衰减，而不是形成垂直传播的光线。这也解释电子为什么通常必须靠近纳米结构才能有效采样同步分量。这里的 $\ell_{\rm ev}$ 是\emph{场幅}衰减长度，强度衰减长度是其一半。

\begin{exbox}[由实场峰值解析求边带与能宽]
给定实场峰值 $E_{\rm pk}^{\rm(real)}$、作用长度 $L_{\rm int}$ 和失配 $\Delta\kappa$，求弱边带权重，并用完整能谱宽度反求场幅。

\textbf{解。} 令
\begin{equation*}
 b=|\beta|=\frac{eE_{\rm pk}^{\rm(real)}L_{\rm int}}{2\hbar\omega}
 \left|\sinc\!\left(\frac{\Delta\kappa L_{\rm int}}2\right)\right|.
\end{equation*}
用 Bessel 幂级数而不逐阶数值求值：
\begin{equation*}
 J_0(2b)=1-b^2+\frac{b^4}{4}+O(b^6),\quad
 J_1(2b)=b-\frac{b^3}{2}+O(b^5),\quad
 J_2(2b)=\frac{b^2}{2}+O(b^4).
\end{equation*}
平方后得到
\begin{equation*}
 P_0=1-2b^2+\frac32b^4+O(b^6),\quad
 P_{\pm1}=b^2-b^4+O(b^6),\quad
 P_{\pm2}=\frac{b^4}{4}+O(b^6).
\end{equation*}
检查归一化时必须把正负阶都算入：$P_0+2P_1+2P_2=1+O(b^6)$。四阶项恰好抵消，是检查展开是否遗漏系数的简便方法。

若不要求 $b\ll1$，仍可使用完整 Bessel 分布的矩恒等式 $\operatorname{Var}\ell=2b^2$。于是交换导致的能量标准差为
\begin{equation*}
 \sigma_{E,\rm ex}=\sqrt2\hbar\omega b
 =\frac{eE_{\rm pk}^{\rm(real)}L_{\rm int}}{\sqrt2}
 \left|\sinc\!\left(\frac{\Delta\kappa L_{\rm int}}2\right)\right|.
\end{equation*}
已知失配时可直接反解 $E_{\rm pk}^{\rm(real)}$。同步处展宽与场强、长度线性相关；接近 $\sinc$ 零点时反演分母很小，场幅误差被放大。这里使用的是相互作用增加的完整能谱方差；源能宽、仪器展宽与有限窗口的处理见\secref{sec:electron-detector-inversion-worked}。
\end{exbox}

\figref{fig:electron-field-calibration} 给出同一窗口在不同失配下的耦合及其边带概率。
\begin{figure}[htbp]
 \centering
 \includegraphics[width=0.92\linewidth]{electron_field_calibration.pdf}
 \caption{从有限长度场积分到电子谱。左图为复振幅相加形成的 $\sinc$ 失配响应；右图取 $\SI{200}{\kilo\electronvolt}$ 电子、$\SI{800}{\nano\metre}$ 光、$\SI{2}{\micro\metre}$ 同步作用区和 $\SI{1}{\mega\volt\per\metre}$ 实场峰值，此时 $|\beta|=0.6452$，分别显示增益和损失阶数。实场峰值与正频振幅的二倍关系已纳入计算。}
 \label{fig:electron-field-calibration}
\end{figure}

\begin{exbox}[高斯近场、双区域相消与完整边带分布]
矩形窗口可以换成平滑局域场。取正频分量
$E_{\omega,z}(z)=E_*\exp[-z^2/(2d^2)]\ee^{\ii k z}$，其中 $d$ 是场幅包络的标准差，
$E_*$ 是复正频振幅。求单区耦合，再将两个这样的局域作用区组合，判断何种相位能使总边带消失。

\textbf{解。} 记 $\Delta k=k-\omega/v_0$，把轨迹积分完成平方：
\begin{align*}
 \int_{\RR}\exp[-z^2/(2d^2)+\ii\Delta k z]\dd z
 &=\exp[-d^2\Delta k^2/2]
 \int_{\RR}\exp[-(z-\ii d^2\Delta k)^2/(2d^2)]\dd z\\
 &=\sqrt{2\pi}d\exp[-d^2\Delta k^2/2].
\end{align*}
最后一步也可由实高斯的 Fourier 变换得到，不要求对复路径积分作近似。于是
\begin{equation}
 \beta_G=\frac{q_{\mathrm e}E_*\sqrt{2\pi}d}{\hbar\omega}
 \exp[-d^2\Delta k^2/2].
 \label{eq:feqo-gaussian-near-field-coupling}
\end{equation}
能谱只确定这个复耦合的模长；场强、长度与失配不能仅凭一个能谱各自唯一反演。

对两区干涉必须说明空间载波的相位约定。取共享载波 $\ee^{\ii kz}$ 的两个包络，中心 $z_{1,2}=\mp D/2$，
后一包络多乘 $\ee^{\ii\chi}$。令 $u=z-z_j$ 后每区的积分多乘 $\ee^{\ii\Delta k z_j}$，故
\begin{equation}
 \beta_{\rm two}=\beta_G(\ee^{-\ii\Delta k D/2}+\ee^{\ii\chi}\ee^{\ii\Delta k D/2}),\qquad
 |\beta_{\rm two}|^2=4|\beta_G|^2\cos^2[(\Delta kD+\chi)/2].
 \label{eq:feqo-two-gaussian-field-interference}
\end{equation}
两区的单独耦合非零时，仍可令 $\chi=\pi-\Delta kD$ 使总耦合为零。
若实际使用各区自身的载波 $\ee^{\ii k(z-z_j)}$，相位参数需相应平移，不能沿用上述 $\chi$ 的定义。
复数相加还要求区间内已占据边带的相对色散相位可忽略；否则须使用
$U_2U_{\rm drift}U_1$，而不能把传播作用藏进一个场强平均。

取 $K_0=\SI{200}{\kilo\electronvolt}$、$\lambda=\SI{800}{\nano\metre}$、$k=0$，
$|E_*|=2\times10^7\,\si{\volt\per\metre}$、$d=\SI{100}{\nano\metre}$。
相对论速度 $v_0=2.08450\times10^8\,\si{\metre\per\second}$ 给
$|\Delta k|d=1.12956$、$|\beta_G|=1.70917$。
这里实场峰值是 $2|E_*|$。弱场近似 $P_{\pm1}\simeq|\beta_G|^2=2.9213$ 已超出概率范围，必须回到
$P_\ell=J_\ell^2(2|\beta_G|)$：$P_{\pm1}\simeq0.02944$、$P_{\pm2}\simeq0.21893$。
单光子边带较弱来自 Bessel 振幅相消，并不意味着相互作用弱。
\end{exbox}

\begin{figure}[htbp]
 \centering
 \includegraphics[width=0.76\linewidth]{pinem_gaussian_example.pdf}
 \caption{上述高斯近场算例的完整边带人口，$|\beta_G|=1.70917$。
 场参数表示正频振幅；边带最大人口在 $\ell=\pm2$，弱场单光子近似已失效。}
 \label{fig:pinem-gaussian-example}
\end{figure}

\begin{figure}[htbp]
 \centering
 \includegraphics[width=0.76\linewidth]{pinem_bessel.pdf}
 \caption{完整 Bessel 谱在 $2|\beta|=5$ 时的分布。正负阶等人口；$\ell=\pm2$ 接近 $J_2$ 的零点而被压低，
 $\ell=\pm4$ 的人口最大。各阶的相消与总概率守恒同时成立。}
 \label{fig:pinem-bessel-sidebands}
\end{figure}


\section{有限脉冲与空间几何}

实际光脉冲有有限持续时间，近场也只分布在样品附近。不同到达时刻、不同横向位置的电子因此经历不同的场积分。本节把这些差别放回 $\beta$ 中，求有限电子波包的输出振幅，并说明什么时候应当相加振幅，什么时候应当平均概率。

\subsection{有限脉冲与相干平均}

若电子波包的持续时间与光脉冲相当，波包前后各部分受到的调制通常不同。前面得到的 $\beta(\zeta,\tau)$ 已保留这种依赖：$\zeta$ 标记共动波包内的位置，$\tau$ 表示电子与光的相对延迟。在程函近似与通道平移对称仍成立时，每个位置对应的散射相位为
\begin{equation}
{
\mathcal U_{\rm p}(\zeta,\tau)
=\exp[\beta(\zeta,\tau)\ee^{\ii\kappa_\omega\zeta}
-\beta^*(\zeta,\tau)\ee^{-\ii\kappa_\omega\zeta}],
\qquad \kappa_\omega=\omega/v_0.}
\label{eq:pinem_finite_pulse_local_phase}
\end{equation}
与连续波极限相比，改变的不是局域幺正结构，而是不同 $\zeta$ 对应不同的复耦合，因此实际电子输出必须把这些切片的相干振幅与入射纵向波包一起传播或积分。

多量子路径也可以在不指定脉冲形状时计数。令 $\ell_+=\max(\ell,0)$、$\ell_-=\max(-\ell,0)$。对净交换 $\ell$ 个量子的所有吸收/发射配对，用非负整数 $j$ 标记额外的往返对数，则局域路径系数为
\begin{equation}
{
C_j^{(\ell)}(\zeta,\tau)
=\frac{\beta(\zeta,\tau)^{j+\ell_+}
[-\beta^*(\zeta,\tau)]^{j+\ell_-}}
{(j+\ell_+)!(j+\ell_-)!}.}
\label{eq:park_Cjn_general}
\end{equation}
因此有限脉冲的普适变化是：同一净边带中的不同 $j$ 路径具有不同的局域权重，最终测得的边带不再只由一个全局复数 $\beta$ 决定。若入射电子纵向振幅记为 $\psi_i(\zeta)$，则对任意具体脉冲都应先在振幅层构造
\[
\psi_f(\zeta,\tau)
=\mathcal U_{\rm p}(\zeta,\tau)\psi_i(\zeta),
\]
再根据探测基计算概率；只有在 $\beta(\zeta,\tau)$ 在电子支持上近似常量时，才恢复单一 Bessel 边带分布。

\begin{exbox}[高斯光脉冲]
取 $\sigma_p>0$，切片耦合为
\[
\beta(\zeta,\tau)=\beta_0
\exp\!\left[-\frac{(\zeta+v_0\tau)^2}{4v_0^2\sigma_p^2}\right].
\]
求第 $j$ 条往返路径的有效包络宽度，并说明有限电子包络怎样进入测得的边带。

\textbf{解。} 第 $j$ 条、净交换数为 $\ell$ 的路径含总次数
$r+s=2j+|\ell|$，所以其包络权重额外带有
\[
\exp\!\left[-\frac{(2j+|\ell|)(\zeta+v_0\tau)^2}
{4v_0^2\sigma_p^2}\right].
\]
因此更高阶往返路径采样更窄的有效时窗。若电子自身也有有限纵向包络，实际边带还要对电子振幅作相干积分；短电子极限恢复局域 $J_\ell^2[2|\beta(\zeta_0,\tau)|]$，而长电子或未锁相平均会逐渐抹去载波相位敏感的交叉项。
\end{exbox}

边带归一化直接来自局域相位算符的幺正性，Bessel 恒等式只是这一幺正性的级数表示。有限脉冲下若 $\beta(\zeta,\tau)$ 随切片变化，不同切片的 $A_\ell$ 具有不同复相位；纵向态保持相干时，应先按传播与探测模式叠加振幅，再取模平方。

\begin{definition}[不同制备事件的 PINEM 系综]
令 $f(\mathbf R,t_0)$ 是电子横向位置和到达时刻的归一化概率密度。
每次事件在自己的轨迹上经历耦合 $\beta(\mathbf R,t_0)$；未记录事件标签时，
输出为相应纯态的经典混合，能谱为
\begin{equation}
\overline P_\ell=\int f(\mathbf R,t_0)
J_\ell^2(2|\beta(\mathbf R,t_0)|)\,\dd^2R\dd t_0.
\label{eq:feqo-ensemble-pinem}
\end{equation}
\label{chap:feqo-averaging}
这不同于同一电子内部包络切片之间的相干叠加；后者必须先求总振幅。
\end{definition}

\begin{exbox}[束腰与光脉冲宽度怎样改变边带人口]
取二维高斯束 $f_R(R)=\exp[-R^2/(2\sigma_R^2)]/(2\pi\sigma_R^2)$，
场积分的横向剖面为 $\beta(R)=\beta_0\exp[-R^2/(2w_\beta^2)]$。
求弱驱动下的平均第一边带概率，再加入独立的高斯到达时刻抖动，比较结果与轴上准时入射。

\textbf{解。} 若 $|\beta_0|\ll1$，由 $J_1(2b)=b+O(b^3)$ 得
\begin{align*}
\overline P_1&=|\beta_0|^2\int_0^\infty
\frac{R\dd R}{\sigma_R^2}
\exp\left[-\left(\frac1{2\sigma_R^2}+\frac1{w_\beta^2}\right)R^2\right]
+O(|\beta_0|^4)\\
&=\frac{|\beta_0|^2}{1+2\sigma_R^2/w_\beta^2}+O(|\beta_0|^4).
\end{align*}
若还独立有高斯到达时刻分布，标准差为 $\sigma_t$，并取
$\beta(t_0)\propto\exp[-t_0^2/(2\tau_L^2)]$，则多出一维积分因子
$[1+2\sigma_t^2/\tau_L^2]^{-1/2}$。
例如 $\sigma_R=w_\beta$、$\sigma_t=\tau_L$ 时，总第一边带人口
是轴上准时电子的 $1/(3\sqrt3)$ 倍，而不是某个平均振幅的平方。
强场时应使用原 Bessel 积分，不能继续沿用这个低阶概率。
\end{exbox}

\begin{proposition}[强度平均与参考相位平均读取不同态分量]
若所有事件有相同 $b=|\beta|$，只有相位 $\phi$ 随事件改变，则
\[
\overline\rho_{\ell m}=J_\ell(2b)J_m(2b)
\langle\ee^{\ii(\ell-m)\phi}\rangle.
\]
所以能谱保持 $J_\ell^2(2b)$。圆周均匀相位使所有非对角元为零；
参考附近标准差为 $\sigma_\phi$ 的高斯相位使其乘
$\exp[-(\ell-m)^2\sigma_\phi^2/2]$。
\end{proposition}
\begin{proof}
逐事件的振幅为 $c_\ell=J_\ell(2b)\ee^{\ii\ell\phi}$。
对 $c_\ell c_m^*$ 作概率平均即得第一式；两个分布的 Fourier 特征函数给出后两式。
因此随机参考相位可以完全抹去后续聚束与相位读出，却不必改变单次能谱人口。
\end{proof}


\subsection{Green 递推与经典耦合}

也可以直接在能量通道中求同一个低反冲散射问题。电子每吸收或发射一个量子，振幅就转移到相邻通道；自由 Green 函数描述两次交换之间的传播。把这些路径逐次相加，可以得到边带振幅的递推关系，并检验它与轨迹相位法给出的 Bessel 解一致。

对中心速度 $v_0$ 的窄带电子，选中的单色近场只连接相差 $\hbar\omega$ 的通道。把电子能量格写成 $|\ell\rangle$，在平移对称极限中相邻边的复交换率相同，时间积分定义的无量纲耦合就是
\begin{equation}
{
\beta
=\frac{q_{\mathrm e}}{\hbar\omega}
\int_{-\infty}^{\infty}\dd z\,
E_{\omega,z}(z)\ee^{-\ii\omega z/v_0}.}
\label{eq:beta_master}
\end{equation}
\eqnrefs{eq:beta_master} 是脉冲耦合~\eqnrefs{eq:pinem_standard_beta_pulsed} 在包络对电子近似恒定时的极限：$|\beta|$ 是累计相干交换强度，$\arg\beta$ 是电子相对于光学场的相位基准。

在相位坐标 $\theta$ 中，经典散射算符写成
\begin{equation*}
S_{\rm cl}(\beta)|\theta\rangle
=\exp[\beta \ee^{-\ii\theta}-\beta^*\ee^{\ii\theta}]|\theta\rangle.
\end{equation*}
在固定参考面，以中心电子通过时刻为 $t=0$，能量阶数 $\ell$ 的时间因子为 $\ee^{-\ii\ell\omega t}$。因此与 $\beta=|\beta|\ee^{\ii\varphi_\beta}$ 一致的时间相位为 $\varphi_{\rm cl}(t)=2|\beta|\sin(\varphi_\beta-\omega t)$，其瞬时能量偏移是相位的负时间梯度：
\begin{equation}
{
\delta E_{\rm inst}(t)
=-\hbar\partial_t\varphi_{\rm cl}(t)
=2\hbar\omega|\beta|\cos(\varphi_\beta-\omega t).}
\label{eq:classical_energy_mod_from_beta}
\end{equation}
而能量本征测量把同一个周期相位 Fourier 分解成离散边带。利用 Bessel 生成函数得到算符展开
\begin{equation}
{
S_{\rm cl}(\beta)
=\sum_{\ell=-\infty}^{\infty}
J_\ell(2|\beta|)\ee^{\ii\ell\varphi_\beta}(\hat b^\dagger)^\ell.}
\label{eq:pinem_fourier_expansion}
\end{equation}
所以“连续能量调制”和“离散 PINEM 边带”是同一散射相位在时间/能量两个共轭表象中的描述。对单能入射和纯相位调制，$J_{-\ell}^2=J_{\ell}^2$ 使增益与损失完全对称，因此平均边带阶数为
\begin{equation}
{\sum_{\ell=-\infty}^{\infty}\ell J_\ell^2(2|\beta|)=0.}
\label{eq:pinem_zero_mean}
\end{equation}
平均能量保持不变并不意味着没有能量展宽：方差随 $|\beta|$ 增大，而平均值因对称性保持零。

单能入射、双边均匀能量格与纯相位散射共同定义理想 Bessel 有效模型 $P_\ell=J_\ell^2(2|\beta|)$。在这一模型内，边带阶数的统计矩可由特征函数统一求得。定义
\begin{equation}
{
\chi_\ell(\vartheta)
\equiv\sum_{\ell=-\infty}^{\infty}P_\ell \ee^{\ii\ell\vartheta}
=J_0\!\left(4|\beta|\sin\frac{\vartheta}{2}\right).}
\label{eq:pinem-sideband-characteristic-dp83}
\end{equation}
在 $\vartheta=0$ 处求导得到边带阶数的偶数矩：
\begin{equation}
{
\langle \ell^2\rangle=2|\beta|^2,
\qquad
\langle \ell^4\rangle=2|\beta|^2+6|\beta|^4.}
\label{eq:pinem-sideband-moments-dp83}
\end{equation}
因此对单能入射，输出能量方差为
\begin{equation}
{
\operatorname{Var}(E_{\rm out})
=2|\beta|^2(\hbar\omega)^2,}
\label{eq:pinem-energy-variance-dp83}
\end{equation}
而不是一个只凭“边带数大约为 $2|\beta|$”估计的经验宽度。四阶矩还给出对非高斯尾部的直接检验；它与二阶矩一起可用于判断实验谱是否仍由单一纯相位 PINEM 核描述。把 $\ell\hbar\omega$ 解释为真实电子输出能量变化还要求均匀格近似在显著占据窗口内成立；一旦真实通道能量 $\mathcal E_\ell$ 的反冲失谐不可忽略，能量矩必须改用 $\mathcal E_\ell$ 加权，而不能继续把\eqnrefs{eq:pinem-energy-variance-dp83} 当作精确电子色散结果。

\begin{lemma}[Bessel 边带分布的特征函数与矩]
设 $\beta\in\CC$，$P_\ell=J_\ell^2(2|\beta|)$，$\ell\in\ZZ$。则其特征函数为 $J_0(4|\beta|\sin(\vartheta/2))$，概率和为一，所有奇数矩为零，$\langle\ell^2\rangle=2|\beta|^2$，$\langle\ell^4\rangle=2|\beta|^2+6|\beta|^4$。仅在均匀能量格中能量方差等于\eqnrefs{eq:pinem-energy-variance-dp83}。
\end{lemma}
\begin{proof}
令
\begin{equation*}
 x\equiv2|\beta|.
\end{equation*}
正文纯相位散射因子具有 Fourier 展开
\begin{equation*}
 \ee^{\ii x\sin\phi}
 =\sum_{\ell=-\infty}^{\infty}
 J_\ell(x)\ee^{\ii\ell\phi}.
\end{equation*}
对于理想单能入射，输出第 $\ell$ 条边带的振幅只差一个整体相位，因此
\begin{equation*}
 P_\ell=J_\ell^2(x).
\end{equation*}
为了同时获得归一化和所有矩，不逐项求和 $\sum_\ell \ell^nJ_\ell^2(x)$，而构造概率特征函数
\begin{equation*}
 \chi_\ell(\vartheta)
 =\sum_{\ell=-\infty}^{\infty}
 P_\ell \ee^{\ii\ell\vartheta}.
\end{equation*}

考虑同一个相位因子的角向自相关
\begin{equation*}
 C(\vartheta)
 \equiv\frac1{2\pi}
 \int_0^{2\pi}\dd\phi\,
 \ee^{\ii x\sin(\phi+\vartheta)}
 \ee^{-\ii x\sin\phi}.
\end{equation*}
先用 Fourier 展开计算。第一因子为
\begin{equation*}
 \ee^{\ii x\sin(\phi+\vartheta)}
 =\sum_{\ell}J_\ell(x)
 \ee^{\ii\ell\phi}\ee^{\ii\ell\vartheta},
\end{equation*}
第二因子可写成
\begin{equation*}
 \ee^{-\ii x\sin\phi}
 =\sum_mJ_m(x)\ee^{-\ii m\phi},
\end{equation*}
这里利用 $J_m(x)$ 对实 $x$ 为实数。于是
\begin{align*}
 C(\vartheta)
 &=\frac1{2\pi}
 \sum_{\ell,m}J_\ell(x)J_m(x)\ee^{\ii\ell\vartheta}
 \int_0^{2\pi}\dd\phi\,
 \ee^{\ii(\ell-m)\phi}\\
 &=\sum_{\ell,m}J_\ell(x)J_m(x)\ee^{\ii\ell\vartheta}
 \delta_{\ell m}\\
 &=\sum_\ell J_\ell^2(x)\ee^{\ii\ell\vartheta}
 =\chi_\ell(\vartheta).
\end{align*}
因此自相关本身就是边带阶数的特征函数。

另一方面，不使用 Fourier 级数也可以直接算同一个积分。三角恒等式给
\begin{align*}
 \sin(\phi+\vartheta)-\sin\phi
 &=2\sin\frac{\vartheta}{2}
 \cos\!\left(\phi+\frac{\vartheta}{2}\right).
\end{align*}
令
\begin{equation*}
 u=\phi+\frac{\vartheta}{2},
\end{equation*}
由于积分区间跨越完整 $2\pi$ 周期，平移不改变积分：
\begin{align*}
 C(\vartheta)
 &=\frac1{2\pi}\int_0^{2\pi}\dd u\,
 \exp\!\left[
 \ii\,2x\sin\frac{\vartheta}{2}\cos u
 \right].
\end{align*}
Bessel 函数的积分定义
\begin{equation*}
 J_0(y)
 =\frac1{2\pi}\int_0^{2\pi}\dd u\,\ee^{\ii y\cos u}
\end{equation*}
于是给出
\begin{align*}
 \chi_\ell(\vartheta)
 &=J_0\!\left(2x\sin\frac{\vartheta}{2}\right)\\
 &=J_0\!\left(4|\beta|\sin\frac{\vartheta}{2}\right),
\end{align*}
即正文\eqnrefs{eq:pinem-sideband-characteristic-dp83}。令 $\vartheta=0$：
\begin{equation*}
 \sum_\ell P_\ell
 =\chi_\ell(0)=J_0(0)=1,
\end{equation*}
所以边带分布的归一化也同时得到，而不是另行假定。

特征函数的 $n$ 阶导数满足
\begin{equation*}
 \chi_\ell^{(n)}(0)
 =\sum_\ell(\ii\ell)^nP_\ell
 =\ii^n\langle\ell^n\rangle.
\end{equation*}
由于 $P_{-\ell}=P_\ell$，所有奇数矩为零。为了得到二阶和四阶矩，先把复合函数在 $\vartheta=0$ 附近展开。令
\begin{equation*}
 y(\vartheta)
 =4|\beta|\sin\frac{\vartheta}{2}.
\end{equation*}
正弦展开为
\begin{align*}
 y(\vartheta)
 &=4|\beta|
 \left(\frac\vartheta2-\frac1{6}\frac{\vartheta^3}{8}
 +O(\vartheta^5)\right)\\
 &=2|\beta|\vartheta
 -\frac{|\beta|}{12}\vartheta^3
 +O(\vartheta^5).
\end{align*}
因此
\begin{align*}
 y^2
 &=4|\beta|^2\vartheta^2
 -\frac{|\beta|^2}{3}\vartheta^4
 +O(\vartheta^6),\\
 y^4
 &=16|\beta|^4\vartheta^4
 +O(\vartheta^6).
\end{align*}
再用
\begin{equation*}
 J_0(y)
 =1-\frac{y^2}{4}+\frac{y^4}{64}+O(y^6),
\end{equation*}
逐项代入：
\begin{align*}
 \chi_\ell(\vartheta)
 &=1-\frac14
 \left(4|\beta|^2\vartheta^2
 -\frac{|\beta|^2}{3}\vartheta^4\right)
 +\frac1{64}
 16|\beta|^4\vartheta^4
 +O(\vartheta^6)\\
 &=1-|\beta|^2\vartheta^2
 +\left(
 \frac{|\beta|^2}{12}
 +\frac{|\beta|^4}{4}
 \right)\vartheta^4
 +O(\vartheta^6).
\end{align*}
另一方面，按概率矩展开
\begin{align*}
 \chi_\ell(\vartheta)
 &=1
 +\ii\langle\ell\rangle\vartheta
 -\frac{\langle\ell^2\rangle}{2}\vartheta^2
 -\frac{\ii\langle\ell^3\rangle}{6}\vartheta^3
 +\frac{\langle\ell^4\rangle}{24}\vartheta^4
 +O(\vartheta^5).
\end{align*}
比较 $\vartheta^2$ 系数：
\begin{equation*}
 -\frac{\langle\ell^2\rangle}{2}=-|\beta|^2
 \quad\Rightarrow\quad
 \langle\ell^2\rangle=2|\beta|^2.
\end{equation*}
比较 $\vartheta^4$ 系数：
\begin{align*}
 \frac{\langle\ell^4\rangle}{24}
 &=\frac{|\beta|^2}{12}+\frac{|\beta|^4}{4},\\
 \langle\ell^4\rangle
 &=2|\beta|^2+6|\beta|^4.
\end{align*}
这得到正文\eqnrefs{eq:pinem-sideband-moments-dp83}。

最后把边带阶数映射成理想均匀能量格。在该有效模型内
\begin{equation*}
 E_{\rm out}=E_0+\ell\hbar\omega.
\end{equation*}
所以
\begin{align*}
 \langle E_{\rm out}\rangle
 &=E_0+\hbar\omega\langle\ell\rangle
 =E_0,\\
 \operatorname{Var}(E_{\rm out})
 &=\langle(E_{\rm out}-E_0)^2\rangle\\
 &=(\hbar\omega)^2\langle\ell^2\rangle\\
 &=2|\beta|^2(\hbar\omega)^2,
\end{align*}
即正文\eqnrefs{eq:pinem-energy-variance-dp83}。若真实反冲使通道能量不再满足 $\mathcal E_\ell=E_0+\ell\hbar\omega$，最后这一步必须改为直接对真实 $\mathcal E_\ell$ 做加权求和；前面的 Bessel 概率矩只描述理想均匀格模型内部的统计。
\end{proof}




\subsection{相干参考与相位可见性}\label{subsec:pinem-coherent-reference-dp31}

复数 $\beta$ 的模和相位都进入电子振幅，而能谱仪只测各振幅的模平方。如果不同的 $\arg\beta$ 给出同一个能谱，单次能谱就无法确定这个相位。为使相位影响概率，可以让未知调制与一个已知参考调制干涉。下面先写出单次测量丢失了什么，再说明参考场怎样提供缺少的信息。

单次理想 PINEM 能谱只含
\begin{equation*}
 P_\ell=J_\ell^2(2|\beta|),
\end{equation*}
因此能量本征测量对 $\arg\beta$ 严格不敏感。这不是反演算法缺陷，而是当前前向模型本身存在相位盲区：把 $\beta$ 乘以任意整体相位都不会改变任何边带概率。要读取复耦合的相位，必须先用一个已知相位参考把相位差转换成可测的强度变化。

设未知相互作用为 $\beta=|\beta|\ee^{\ii\varphi_\beta}$，随后施加与其锁相的已知参考相互作用 $\beta_R=|\beta_R|\ee^{\ii\varphi_R}$。在双边平移近似下，两个散射算符具有同一生成元，因此精确合成为
\begin{equation}
 {
 S_{\rm cl}(\beta_R)S_{\rm cl}(\beta)
 =S_{\rm cl}(\beta_{\rm tot}),
 \qquad
 \beta_{\rm tot}=\beta+\beta_R.}
 \label{eq:pinem-reference-composition-dp31}
\end{equation}
若两次相互作用之间的传播仅含格点编号的线性相位 $D=\exp(-\ii\varphi_{\rm prop}\Lambda_e)$，则 $D^\dagger b^\dagger D=\ee^{\ii\varphi_{\rm prop}}b^\dagger$，从而 $S_{\rm cl}(\beta_R)DS_{\rm cl}(\beta)=D S_{\rm cl}(\beta+\beta_R\ee^{\ii\varphi_{\rm prop}})$。左侧最终的 $D$ 不改变能谱，参考相位便可写成 $\varphi_R\to\varphi_R+\varphi_{\rm prop}$。二次及更高阶色散相位不能只靠这个相位替换处理，须保留完整的 $S_{\rm cl}(\beta_R)U_{\rm drift}S_{\rm cl}(\beta)$。以下只考虑可合成的情形，定义信号耦合与总参考相位之差 $\delta\equiv\varphi_\beta-\varphi_R$，则
\begin{equation*}
 B^2\equiv|\beta_{\rm tot}|^2
 =|\beta|^2+|\beta_R|^2
 +2|\beta||\beta_R|\cos\delta,
\end{equation*}
最终边带概率变成
\begin{equation}
 {P_\ell(\delta)=J_\ell^2(2B).}
 \label{eq:pinem-reference-sidebands-dp31}
\end{equation}
相位已经通过复振幅相干叠加进入 $B$，因此扫描参考相位就能把 $\arg\beta$ 转换成可观测的边带重分布。

对于理想的完整边带读出，全部边带关于 $\delta$ 的经典 Fisher 信息可解析求和：
\begin{equation}
 {
 F_\delta^{\rm(full)}
 =8\,
 \frac{|\beta|^2|\beta_R|^2\sin^2\delta}
 {|\beta+\beta_R|^2}.}
 \label{eq:pinem-reference-phase-fisher-dp31}
\end{equation}
相位以弧度计，因而 $F_\delta$ 无量纲。令 $r=|\beta|$、$R=|\beta_R|$。对固定的 $r$ 与 $R$，理想完整边带读出的最大局部相位信息为
\begin{equation}
 {
 F_{\delta,\max}^{\rm(full)}=8\min(r^2,R^2),
 \qquad
 \cos\delta_{\rm opt}
 =-\frac{\min(r,R)}{\max(r,R)}.}
 \label{eq:pinem-reference-phase-fisher-optimum-dp32}
\end{equation}
当 $r=R$ 时最优点逼近完全相消方向；这是一个非正则极限，应从边带概率本身连续取极限，而不能把含 $1/B$ 的中间表达式直接代入 $B=0$。正交参考 $\delta=\pi/2$ 更远离这一奇异参数化点，且在 $R=r$ 时仍有 $F_\delta=4r^2$，因此常是稳定而非绝对最优的实验工作点。若参考显著强于未知信号，正交工作点的信息趋近 $8r^2$。有限边带窗口、探测效率和能量展宽都属于对理想边带分布的经典粗粒化，因此只能降低实际可用的 Fisher 信息。

\subsection{反冲修正}

Bessel 解要求相邻边的完整耦合矩阵元近似相同，并且作用包络不能重新解析出被平均掉的快速支路。设 $\mathcal V_N$ 是实际第 $N\to N+1$ 条边的完整矩阵元，在显著占据窗口 $|N|\le N_{\rm occ}$ 内定义
\begin{equation}
{
\epsilon_{\rm eq}
=\max_{|N|\le N_{\rm occ}}
\left|\frac{\mathcal V_N}{\mathcal V_0}-1\right|\ll1.}
\label{eq:feist2015_equal_rabi_recoil_condition}
\end{equation}
这个定义只比较实际矩阵元，不预先把其变化归因于速度、场谱或自旋器重叠。它是瞬时相对误差；在共同失谐且所有相关边均受控的经典双边格上，若共同耦合率为 $\mathcal G_0(t)$，加权平移及其厄米共轭给 $\|H-H_0\|/\hbar\le2\epsilon_{\rm eq}|\mathcal G_0(t)|$，故 Duhamel 界为 $\|U-U_0\|\le2\epsilon_{\rm eq}\int|\mathcal G_0(t)|\dd t$。因此多次强交换还要检查累计误差，实际工作窗外的通道需另估计。

为了把\eqnrefs{eq:feist2015_equal_rabi_recoil_condition} 分解成可计算来源，随后代入一维纵向、$\phi=0$、固定自旋的正能 Dirac 通道。先在长度 $L_z$ 的周期盒中取
\begin{equation*}
 \langle z|p_N,s,+\rangle
 =\frac{1}{\sqrt{L_z}}\ee^{\ii p_Nz/\hbar}u_s(p_N),
 \qquad u_s^\dagger(p_N)u_s(p_N)=1,
\end{equation*}
并令 $\Delta\kappa_N\equiv(p_{N+1}-p_N)/\hbar$。若纵向矢势可分离为 $A_z(z,t)=a(t)A_z(z)$，定义盒归一 Fourier 分量
\begin{equation*}
 \widetilde A_z(\kappa)
 \equiv\frac1{L_z}\int_0^{L_z}\dd z\,A_z(z)\ee^{-\ii\kappa z}
\end{equation*}
以及正能 Dirac 电流矩阵元
\begin{equation*}
 \mathcal J_N
 \equiv c\,u_s^\dagger(p_{N+1})\alpha_z u_s(p_N).
\end{equation*}
把盒归一态代入 $\hat V_A=-q_{\mathrm e} c\alpha_zA_z$，逐项积分得到
\begin{align*}
\mathcal V_N
&\equiv {}_V\!\langle p_{N+1},s,+|\hat V_A|p_N,s,+\rangle_V\\
&=-\frac{q_{\mathrm e} c}{L_z}
\int_0^{L_z}\dd z\,
 u_s^\dagger(p_{N+1})\alpha_z u_s(p_N)
 A_z(z)\ee^{-\ii(p_{N+1}-p_N)z/\hbar}\\
&=-q_{\mathrm e}\,\mathcal J_N\,\widetilde A_z(\Delta\kappa_N).
\end{align*}
因此
\begin{equation}
 {
 \mathcal V_N
 =-q_{\mathrm e}\,\mathcal J_N\,\widetilde A_z(\Delta\kappa_N).}
 \label{eq:pinem-edge-matrix-element-dirac-dp116}
\end{equation}
于是定义
\begin{equation*}
 \epsilon_J
 \equiv\max_N\left|\frac{\mathcal J_N}{\mathcal J_0}-1\right|,
 \qquad
 \epsilon_A
 \equiv\max_N\left|
 \frac{\widetilde A_z(\Delta\kappa_N)}
 {\widetilde A_z(\Delta\kappa_0)}-1\right|,
\end{equation*}
便得到严格误差界
\begin{equation}
 {
 \epsilon_{\rm eq}
 \le\epsilon_J+\epsilon_A+\epsilon_J\epsilon_A.}
 \label{eq:pinem-equal-coupling-error-dp68}
\end{equation}
因此逐边变化有两个独立来源：电子正能电流矩阵元的变化和场的纵向空间谱变化。只有再取同自旋、小动量步长极限，才有
\begin{equation*}
 \lim_{p_{N+1}\to p_N}\mathcal J_N
 =c\,u_s^\dagger(p_N)\alpha_z u_s(p_N)
 =\frac{\partial E}{\partial p}\bigg|_{p_N}
 =v_N,
\end{equation*}
此时 $\epsilon_J$ 才进一步约化为常见的速度色散判据。

\begin{proposition}[逐边矩阵元的乘积相对误差]
设中心矩阵元 $\mathcal V_0$ 及两因子 $\mathcal J_0,\widetilde A_z(\Delta\kappa_0)$ 均非零。Dirac 电流与场 Fourier 系数因子化后，逐边相对误差满足 $|\mathcal V_N/\mathcal V_0-1|\le|\delta_J^{(N)}|+|\delta_A^{(N)}|+|\delta_J^{(N)}\delta_A^{(N)}|$，从而得到\eqnrefs{eq:pinem-equal-coupling-error-dp68}。中心场为零或相消节点附近不能用此相对误差判据。
\end{proposition}
\begin{proof}
在选定的一维纵向模型中，线性相互作用算符为
\begin{equation*}
 \hat V_A(t)=-q_{\mathrm e} c\alpha_z a(t)A_z(\hat z).
\end{equation*}
周期盒正能态满足
\begin{equation*}
 \langle z|p_N,s,+\rangle
 =L_z^{-1/2}\ee^{\ii p_Nz/\hbar}u_s(p_N).
\end{equation*}
因此去掉共同时间包络 $a(t)$ 后，第 $N\to N+1$ 条边的矩阵元为
\begin{align*}
 \mathcal V_N
 &=\langle p_{N+1},s,+|
 (-q_{\mathrm e} c\alpha_zA_z)
 |p_N,s,+\rangle\\
 &=-\frac{q_{\mathrm e} c}{L_z}
 u_s^\dagger(p_{N+1})\alpha_z u_s(p_N)
 \int_0^{L_z}\dd z\,
 A_z(z)\ee^{-\ii(p_{N+1}-p_N)z/\hbar}\\
 &=-q_{\mathrm e}\,\mathcal J_N\,
 \widetilde A_z(\Delta\kappa_N),
\end{align*}
其中
\begin{equation*}
 \mathcal J_N=c\,u_s^\dagger(p_{N+1})\alpha_z u_s(p_N),
 \qquad
 \Delta\kappa_N=\frac{p_{N+1}-p_N}{\hbar}.
\end{equation*}
这得到正文\eqnrefs{eq:pinem-edge-matrix-element-dirac-dp116}。

假定 $\mathcal V_0\neq0$、$\mathcal J_0\neq0$ 且 $\widetilde A_z(\Delta\kappa_0)\neq0$，则
\begin{align*}
 \frac{\mathcal V_N}{\mathcal V_0}
 &=\frac{\mathcal J_N}{\mathcal J_0}
 \frac{\widetilde A_z(\Delta\kappa_N)}
 {\widetilde A_z(\Delta\kappa_0)}\\
 &=\left[1+\delta_J^{(N)}\right]
 \left[1+\delta_A^{(N)}\right]\\
 &=1+\delta_J^{(N)}+\delta_A^{(N)}
 +\delta_J^{(N)}\delta_A^{(N)},
\end{align*}
其中
\begin{equation*}
 \delta_J^{(N)}\equiv\frac{\mathcal J_N}{\mathcal J_0}-1,
 \qquad
 \delta_A^{(N)}\equiv
 \frac{\widetilde A_z(\Delta\kappa_N)}
 {\widetilde A_z(\Delta\kappa_0)}-1.
\end{equation*}
于是对 $|N|\le N_{\rm occ}$ 取最大值并使用三角不等式：
\begin{align*}
 \epsilon_{\rm eq}
 &=\max_N\left|
 \delta_J^{(N)}+\delta_A^{(N)}
 +\delta_J^{(N)}\delta_A^{(N)}\right|\\
 &\le\epsilon_J+\epsilon_A+\epsilon_J\epsilon_A,
\end{align*}
这就是\eqnrefs{eq:pinem-equal-coupling-error-dp68}。

最后说明速度因子的来源。对同一归一化正能本征态，Hellmann--Feynman 恒等式给
\begin{align*}
 \frac{\partial E(p)}{\partial p}
 &=u_s^\dagger(p)
 \frac{\partial}{\partial p}
 \left(c\alpha_z p+\beta_Dm_{\mathrm e} c^2\right)
 u_s(p)\\
 &=c\,u_s^\dagger(p)\alpha_z u_s(p).
\end{align*}
在 $p_{N+1}\to p_N$ 的小步长极限，$\mathcal J_N\to v_N$。
\end{proof}


对恰好同步、快慢支路具有相同常数矩阵元的模型，设开关包络 $w(t)$ 的积分非零，并选择归一化 $\widetilde w(0)=1$。这里 $w$ 的单位为 $\mathrm{s^{-1}}$，实际交换强度已提取到它前面的无量纲系数中。共振支路在旋转表象中是慢项，而被舍弃的反旋支路相对它多出 $2\omega$ 的快速相位；剥去共同矩阵元后，两条支路的一阶时间积分之比为
\begin{equation}
{
\frac{\mathcal A_{\rm fast}^{(1)}}{\mathcal A_{\rm slow}^{(1)}}
=\frac{\widetilde w(2\omega)}{\widetilde w(0)}
=\widetilde w(2\omega),
\qquad
\widetilde w(\Omega)=\int\dd t\,w(t)\ee^{\ii\Omega t}.}
\label{eq:pinem_fast_branch_envelope_ratio}
\end{equation}
因此一阶反旋跃迁相对于同步跃迁的抑制可用
\begin{equation}
{
\epsilon_{2\omega}=|\widetilde w(2\omega)|\ll1.}
\label{eq:feist2015_long_interaction_condition}
\end{equation}
对矩形时间窗它退化成 sinc；对高斯脉冲、啁啾脉冲或脉冲列必须使用各自的 $\widetilde w$。若慢支路的共同失谐为 $\Delta_0$，应比较 $\widetilde w(\Delta_0+2\omega)$ 与 $\widetilde w(\Delta_0)$；若后者为零，就改用各支路的绝对振幅，不能继续定义这个比值。快慢矩阵元不同还需乘上它们的比值。一阶积分很小只控制真实跃迁的最低阶振幅，反旋往返产生的相位仍可随作用时间累积。

要控制整个传播子，需要同时处理真实跃迁与反旋虚过程。把已选工作空间中的哈密顿量除以 $\hbar$，写成
\[
 h(t)=h_s(t)+F(t),\qquad
 F(t)=K(t)\ee^{\ii\Sigma t}+K^\dagger(t)\ee^{-\ii\Sigma t},
\]
其中 $h_s=h_s^\dagger$ 是保留的慢生成元，$K$ 是被舍支路的算符系数，二者的单位均为 $\mathrm{s^{-1}}$，$\Sigma\ne0$ 是快速失谐。令 $U$ 与 $U_s$ 分别由 $h$ 与 $h_s$ 生成。下面的界适用于有界算符；有限通道模型若由更大的电子与光子空间截断而来，其截断传播误差需另外控制，不能仅因初态光子数有限就使用全程有限算符范数。

\begin{proposition}[快速支路的有限时间传播界]
设 $h_s(t)$ 自伴且有界，$K(t)$ 在算符范数中绝对连续，相关范数在 $[t_i,t_f]$ 上可积。记 $k(t)=\|K(t)\|$。则
\begin{equation}
\begin{aligned}
 \|U(t_f,t_i)-U_s(t_f,t_i)\|
 \le{}&\frac{2[k(t_i)+k(t_f)]}{|\Sigma|}
 +\frac{2}{|\Sigma|}\int_{t_i}^{t_f}\|\dot K(t)\|\,\dd t\\
 &+\frac{4}{|\Sigma|}\int_{t_i}^{t_f}
 \bigl[\|h_s(t)\|k(t)+k(t)^2\bigr]\dd t .
\end{aligned}
\label{eq:feqo-fast-branch-propagator-bound}
\end{equation}
这是一个充分误差界，不要求把 Dyson 级数截到某一阶。
\end{proposition}
\begin{proof}
为保留振荡造成的抵消，先构造无量纲的自伴算符
\[
 C(t)=\frac{K(t)\ee^{\ii\Sigma t}}{\ii\Sigma}
       -\frac{K^\dagger(t)\ee^{-\ii\Sigma t}}{\ii\Sigma}.
\]
对它求导，快速相位给出 $F$，包络导数给出余项 $D$：
\[
 \dot C=F+D,\qquad
 D=\frac{\dot K\ee^{\ii\Sigma t}
          -\dot K^\dagger\ee^{-\ii\Sigma t}}{\ii\Sigma}.
\]
因此 $\|C\|\le2k/|\Sigma|$、$\|D\|\le2\|\dot K\|/|\Sigma|$，并且 $\|F\|\le2k$。
由两种传播子的 Duhamel 恒等式，
\[
 U-U_s=-\ii\int_{t_i}^{t_f}U(t_f,t)F(t)U_s(t,t_i)\dd t.
\]
把 $F=\dot C-D$ 代入，再对含 $\dot C$ 的项分部积分。传播子对中间时刻的导数分别满足
\[
 \partial_tU(t_f,t)=\ii U(t_f,t)h(t),\qquad
 \partial_tU_s(t,t_i)=-\ii h_s(t)U_s(t,t_i),
\]
从而
\begin{align*}
 \int U\dot C U_s\dd t
 &=[UCU_s]_{t_i}^{t_f}
   -\ii\int U(hC-Ch_s)U_s\dd t,\\
 hC-Ch_s&=[h_s,C]+FC.
\end{align*}
两个传播子均幺正，所以端点项只留下 $\|C(t_i)\|+\|C(t_f)\|$，积分项至多为
\[
 \int\bigl[2\|h_s\|\|C\|+\|F\|\|C\|+\|D\|\bigr]\dd t.
\]
代入三个范数界便得到\eqnrefs{eq:feqo-fast-branch-propagator-bound}。
\end{proof}

这个结果把三种误差分开：脉冲端点带来 $k/|\Sigma|$，包络变化带来 $\int\|\dot K\|/|\Sigma|$，传播中的快慢交叉作用和反旋往返带来累计项。即使端点平滑、一阶 Fourier 积分接近零，仍要检查最后一项。对大致恒定的速率 $k$ 和相近量级的慢生成元，长时控制包含 $k^2T/|\Sigma|\ll1$，不能用 $k/|\Sigma|\ll1$ 单独替代。

\begin{exbox}[一阶快振荡相消以后，剩下什么]
考虑两个归一化通道 $|+\rangle,|-\rangle$，令 $\sigma_z|\pm\rangle=\pm|\pm\rangle$，$\sigma_+=|+\rangle\langle-|$。只保留一个快速支路，取常数 $\gamma>0$、$\Sigma>0$：
\[
 h(t)=\gamma(\sigma_+\ee^{\ii\Sigma t}
                   +\sigma_-\ee^{-\ii\Sigma t}),\qquad 0\le t\le T.
\]
当 $\Sigma T=2\pi m$，$m$ 为正整数，一阶跃迁积分严格为零。求精确传播，并判断能否把这一支路完全删掉。

\textbf{解。} 令 $R(t)=\exp(\ii\Sigma t\sigma_z/2)$，写 $\psi=R\chi$。由 $h=R\gamma\sigma_xR^\dagger$，旋转后的方程为
\[
 \ii\dot\chi=(\gamma\sigma_x+\Sigma\sigma_z/2)\chi.
\]
因为 $\sigma_x$ 与 $\sigma_z$ 反对易，括号的平方为 $\nu^2 I$，其中 $\nu=\sqrt{\gamma^2+\Sigma^2/4}$。指数遂为
\[
 U(t)=R(t)\left[\cos(\nu t)I
 -\ii\frac{\sin(\nu t)}\nu
 (\gamma\sigma_x+\Sigma\sigma_z/2)\right].
\]
对入射通道 $|-\rangle$，整周期端点 $R(T)=(-1)^mI$ 给
\[
 \langle-|U(T)|-\rangle
 =(-1)^m\left[\cos(\nu T)
       +\ii\frac{\Sigma}{2\nu}\sin(\nu T)\right],\qquad
 P_{-\to+}=\frac{\gamma^2}{\nu^2}\sin^2(\nu T).
\]
在 $\gamma/\Sigma\ll1$ 时，$\nu=\Sigma/2+\gamma^2/\Sigma+O(\gamma^4/\Sigma^3)$。若同时 $\gamma^4T/\Sigma^3\ll1$，存活振幅含有累计相位 $\gamma^2T/\Sigma$，而真实转移概率始终至多 $4\gamma^2/\Sigma^2$。所以低末态泄漏不等于传播子接近恒等；入射叠加态可通过两通道的相对相位读出这个虚过程。
\end{exbox}


若等耦合条件失败，最自然的修正不是“改 Bessel 参数”，而是保留自由色散曲率。纵向相对论色散在 $p_0$ 处满足 $E''(p_0)=1/m_\parallel$，所以对 $p_\ell=p_0+\ell\hbar \kappa_z$，
\begin{equation}
{
\mathcal E_\ell
=E_0+\ell\hbar v_0\kappa_z
+\ell^2\hbar\omega_{\rm r}
+O\!\left[(\ell\hbar \kappa_z)^3E^{(3)}(p_0)\right],
\qquad
\omega_{\rm r}=\frac{\hbar \kappa_z^2}{2m_\parallel}.}
\label{eq:pinem_recoil_lattice_expansion}
\end{equation}
若同步条件取 $v_0\kappa_z=\omega$，第 $\ell\to\ell+1$ 条边相对于理想 $\hbar\omega$ 梯的失谐就是
\begin{equation}
{
\Delta_\ell
=\frac{\mathcal E_{\ell+1}-\mathcal E_\ell}{\hbar}-\omega
=(2\ell+1)\omega_{\rm r}+O[(\ell\hbar \kappa_z)^2\kappa_zE^{(3)}(p_0)].}
\label{eq:pinem_recoil_edge_detuning}
\end{equation}
因此占据窗 $|\ell|\le N_{\rm occ}$ 内，忽略二阶反冲至少要求
\begin{equation}
{
\epsilon_{\rm rec}^{(T)}
\equiv(2N_{\rm occ}+1)\omega_{\rm r}T\ll1.}
\label{eq:pinem_recoil_phase_condition}
\end{equation}
若需要把三阶及更高色散的截断误差也纳入同一个有限时间判据，定义
\[
\delta_\ell^{(3)}\equiv
\frac{R_{\ell+1}^{(3)}-R_\ell^{(3)}}{\hbar},
\qquad
\epsilon_{\rm rec}^{(3)}\equiv
T\max_{|\ell|\le N_{\rm occ}}|\delta_\ell^{(3)}|.
\]
逐边总失谐的累计相位满足
\begin{equation}
{
T\max_{|\ell|\le N_{\rm occ}}|\Delta_\ell|
\le
\epsilon_{\rm rec}^{(T)}+\epsilon_{\rm rec}^{(3)}.}
\label{eq:pinem_recoil_total_phase_bound}
\end{equation}
因此使用均匀能量梯时，不仅要检查二次曲率产生的 $\epsilon_{\rm rec}^{(T)}$，还要确认实际占据窗口中的高阶 Taylor 余项没有重新积累出 $O(1)$ 相位。
在非零常幅作用段内，若还要把失谐相对于交换速率本身忽略，可比较 $\epsilon_{\rm rec}^{(\Omega)}=(2N_{\rm occ}+1)\omega_{\rm r}/|\Omega|$；这里 $\Omega$ 是该段的共同交换速率。这个比值适合判断持续交换是否受逐边失谐影响，不是所有短时脉冲的必要条件。短弱脉冲可以有失谐大于耦合率而累计误差很小，强长脉冲则需把逐边差异乘上累计交换强度。两种尺度各有用途，实际 Bessel 精度最终由传播误差控制。\chapref{chap:periodic-recoil}将从\eqnrefs{eq:pinem_recoil_lattice_expansion}重新求有限反冲动力学。

\begin{proposition}[格点色散差的反冲相位界]
设真实通道能量展开为 $\mathcal E_\ell=E_0+\ell\hbar\omega+\ell^2\hbar\omega_{\rm r}+R_\ell^{(3)}$，占据窗 $|\ell|\le N_{\rm occ}$。则逐边失谐为 $(2\ell+1)\omega_{\rm r}+(R_{\ell+1}^{(3)}-R_\ell^{(3)})/\hbar$，其累计模满足\eqnrefs{eq:pinem_recoil_total_phase_bound}；一般实反冲系数需用 $|\omega_{\rm r}|$。
\end{proposition}
\begin{proof}
正文\eqnrefs{eq:pinem_recoil_lattice_expansion} 给出
\[
\mathcal E_\ell
=E_0+\ell\hbar\omega+\ell^2\hbar\omega_{\rm r}+R_\ell^{(3)}.
\]
相邻两条真实电子通道之间的角频率为
\[
\omega_{\ell+1,\ell}
=\frac{\mathcal E_{\ell+1}-\mathcal E_\ell}{\hbar},
\]
故相对于同步频率 $\omega$ 的逐边失谐为
\[
\Delta_\ell
=\frac{\mathcal E_{\ell+1}-\mathcal E_\ell}{\hbar}-\omega.
\]
对离散能量差逐项计算：
\begin{align*}
\mathcal E_{\ell+1}-\mathcal E_\ell
={}&[(\ell+1)-\ell]\hbar\omega
+[(\ell+1)^2-\ell^2]\hbar\omega_{\rm r}
+R_{\ell+1}^{(3)}-R_\ell^{(3)}\\
={}&\hbar\omega+(2\ell+1)\hbar\omega_{\rm r}
+R_{\ell+1}^{(3)}-R_\ell^{(3)}.
\end{align*}
因此
\[
\Delta_\ell
=(2\ell+1)\omega_{\rm r}+\delta_\ell^{(3)},
\qquad
\delta_\ell^{(3)}
=\frac{R_{\ell+1}^{(3)}-R_\ell^{(3)}}{\hbar}.
\]
在显著占据窗口 $|\ell|\le N_{\rm occ}$ 内，二次色散贡献的最大累计相位为
\begin{align*}
T\max_{|\ell|\le N_{\rm occ}}
|(2\ell+1)\omega_{\rm r}|
&=T\omega_{\rm r}
\max_{|\ell|\le N_{\rm occ}}|2\ell+1|\\
&=(2N_{\rm occ}+1)\omega_{\rm r}T\\
&=\epsilon_{\rm rec}^{(T)},
\end{align*}
得到正文\eqnrefs{eq:pinem_recoil_phase_condition}。对完整逐边失谐使用三角不等式：
\begin{align*}
T\max_{|\ell|\le N_{\rm occ}}|\Delta_\ell|
&\le
T\max_{|\ell|\le N_{\rm occ}}|(2\ell+1)\omega_{\rm r}|
+T\max_{|\ell|\le N_{\rm occ}}|\delta_\ell^{(3)}|\\
&=\epsilon_{\rm rec}^{(T)}+\epsilon_{\rm rec}^{(3)},
\end{align*}
即正文\eqnrefs{eq:pinem_recoil_total_phase_bound}。
\end{proof}


速度色散是正能 Dirac 电流矩阵元的低反冲极限，而\eqnrefs{eq:pinem-edge-matrix-element-dirac-dp116} 保留有限步长下的自旋器重叠。

\begin{derivation}[有限脉冲散射的侧带振幅]
固定一个电子切片 $(\zeta,\tau)$，把 $\beta(\zeta,\tau)$ 暂记为 $\beta$，并令
\[
 \theta\equiv\kappa_\omega\zeta.
\]
局域散射算符是纯相位
\begin{equation*}
 U(\theta)=\exp\!\left(\beta \ee^{\ii\theta}-\beta^*\ee^{-\ii\theta}\right).
\end{equation*}
由于两个指数中的量只是复数，它们彼此对易，所以可精确分解为
\begin{equation*}
 U(\theta)
 =\ee^{\beta \ee^{\ii\theta}}
  \ee^{-\beta^*\ee^{-\ii\theta}}.
\end{equation*}
分别作 Taylor 展开，得到双重级数
\begin{align*}
 U(\theta)
 &=\left[\sum_{r=0}^{\infty}
 \frac{\beta^r}{r!}\ee^{\ii r\theta}\right]
 \left[\sum_{s=0}^{\infty}
 \frac{(-\beta^*)^s}{s!}\ee^{-\ii s\theta}\right]\\
 &=\sum_{r,s\ge0}
 \frac{\beta^r(-\beta^*)^s}{r!s!}
 \ee^{\ii(r-s)\theta}.
\end{align*}
这里 $r$ 与 $s$ 分别计数从两个指数中取出的正、负频因子数；真正由能量测量分辨的是它们的差
\[
 \ell=r-s,
\]
而不是 $r$、$s$ 各自。于是同一个最终边带 $\ell$ 包含无穷多条具有相同净交换量、但往返次数不同的振幅路径。

若 $\ell\ge0$，所有非负整数解为
\[
 r=j+\ell,\qquad s=j,
 \qquad j=0,1,2,\ldots;
\]
若 $\ell<0$，则
\[
 r=j,\qquad s=j-\ell.
\]
引入
\[
 \ell_+=\max(\ell,0),
 \qquad
 \ell_-=\max(-\ell,0),
\]
两种情况统一为
\[
 r=j+\ell_+,
 \qquad
 s=j+\ell_-.
\]
因此第 $j$ 条路径对净边带 $\ell$ 的复振幅贡献为
\begin{equation*}
 C_j^{(\ell)}
 =\frac{\beta^{j+\ell_+}(-\beta^*)^{j+\ell_-}}
 {(j+\ell_+)!(j+\ell_-)!},
\end{equation*}
恢复 $(\zeta,\tau)$ 依赖即得到正文\eqnrefs{eq:park_Cjn_general}。这里 $C_j^{(\ell)}$ 是振幅而不是概率；所有具有相同 $\ell$ 的路径必须先相干相加：
\begin{equation*}
 A_\ell(\zeta,\tau)
 =\sum_{j=0}^{\infty}C_j^{(\ell)}(\zeta,\tau).
\end{equation*}

这一求和可以显式完成。写
\[
 \beta=|\beta|\ee^{\ii\varphi_\beta}.
\]
对 $\ell\ge0$，有
\begin{align*}
 A_\ell
 &=\beta^\ell
 \sum_{j=0}^{\infty}
 \frac{(-|\beta|^2)^j}{j!(j+\ell)!}\\
 &=\ee^{\ii\ell\varphi_\beta}
 \sum_{j=0}^{\infty}
 \frac{(-1)^j|\beta|^{2j+\ell}}{j!(j+\ell)!}\\
 &=\ee^{\ii\ell\varphi_\beta}J_\ell(2|\beta|),
\end{align*}
其中最后一步使用 Bessel 函数的级数定义
\[
 J_\ell(2x)
 =\sum_{j=0}^{\infty}
 \frac{(-1)^j x^{2j+\ell}}{j!(j+\ell)!}.
\]
对 $\ell<0$ 使用 $J_{-n}(x)=(-1)^nJ_n(x)$ 得到同一统一结果
\begin{equation*}
 \boxed{
 A_\ell=\ee^{\ii\ell\varphi_\beta}J_\ell(2|\beta|).}
\end{equation*}
因此 Bessel 边带并不是额外假设，而是同一净交换通道内所有往返路径振幅的精确重求和。

还可以直接检验概率守恒。由于指数
$\beta \ee^{\ii\theta}-\beta^*\ee^{-\ii\theta}$
纯虚，任意 $\theta$ 都有 $|U(\theta)|=1$。Fourier 系数定义为
\begin{equation*}
 A_\ell=\frac{1}{2\pi}
 \int_0^{2\pi}\dd\theta\,
 U(\theta)\ee^{-\ii\ell\theta}.
\end{equation*}
Parseval 恒等式于是给
\begin{align*}
 \sum_{\ell\in\mathbb Z}|A_\ell|^2
 &=\frac{1}{2\pi}\int_0^{2\pi}\dd\theta\,|U(\theta)|^2\\
 &=1.
\end{align*}
\end{derivation}

\begin{derivation}[快分支振荡与长作用时间条件]
把归一化开关包络取在有限区间 $[t_i,t_f]$，满足
\[
 \widetilde w(0)=\int_{t_i}^{t_f}\dd t\,w(t)=1,
 \qquad
 \widetilde w(\Omega)=\int_{t_i}^{t_f}\dd t\,w(t)\ee^{\ii\Omega t}.
\]
正文的一阶微扰已经把慢支路与快支路之比化为
$|\widetilde w(2\omega)|$。为了看清“作用时间长”为什么能够压低它，假设 $w$ 在作用区间绝对连续。对 $\Omega\ne0$ 分部积分：
\begin{align*}
 \widetilde w(\Omega)
 &=\left.\frac{w(t)\ee^{\ii\Omega t}}{\ii\Omega}\right|_{t_i}^{t_f}
 -\frac1{\ii\Omega}\int_{t_i}^{t_f}\dd t\,
 \dot w(t)\ee^{\ii\Omega t}.
\end{align*}
取模并用三角不等式得到严格上界
\begin{align*}
 |\widetilde w(\Omega)|
 \le\frac{|w(t_f)|+|w(t_i)|
 +\int_{t_i}^{t_f}\dd t\,|\dot w(t)|}{|\Omega|}.
\end{align*}
因此在 $\Omega=2\omega$ 处，若包络端点平滑降到零，便有
\begin{equation*}
 \left|\frac{\mathcal A_{\rm fast}^{(1)}}{\mathcal A_{\rm slow}^{(1)}}\right|
 =|\widetilde w(2\omega)|
 \le\frac{1}{2\omega}
 \int_{t_i}^{t_f}\dd t\,|\dot w(t)|.
\end{equation*}
这个式子比单纯说“$T$ 很长”更精确：快支路由包络在光学频率尺度上的变化率控制；同样持续时间的两个脉冲，只要边缘平滑程度不同，反旋泄漏就可以不同。

矩形窗口不满足端点为零的平滑条件，需要直接积分。对
$w(t)=T^{-1}\mathbf 1_{[-T/2,T/2]}(t)$，
\begin{align*}
 \widetilde w(2\omega)
 &=\frac1T\int_{-T/2}^{T/2}\dd t\,\ee^{2\ii\omega t}\\
 &=\frac{\sin(\omega T)}{\omega T}
 =\sinc(\omega T).
\end{align*}
所以正文\eqnrefs{eq:feist2015_long_interaction_condition} 的一般小参数是
$\epsilon_{2\omega}=|\widetilde w(2\omega)|$；只有指定矩形开关之后，它才退化为由 $\omega T$ 单独控制的 sinc 条件。
\end{derivation}

